\section{Type \texorpdfstring{$BC'_1$}{BC'1}} \label{section:TypeBC1}

In this section, we deal with the $BC'_1$ case defined in Definition \ref{def:typeBC}.
The aim of this part is to prove that, if $(\lie{g}, \lie{h})$ is of type $BC'_1$, then the tuple $(\lie{g}, \lie{h}, \orbit{O})$ is never minimal.
If $\lie{h}$ is isomorphic to $\lieSp(1,k)$ or $\lieF_{4(-20)}$, then non-minimality is readily verified by simple arguments in Subsection \ref{subsect:EmbeddingBasic}.

In the case where $\lie{h} \simeq \lieSU(1,k)$, there are finitely many possible series for $\lie{g}$, but many candidates occur, especially among exceptional Lie algebras.
For exceptional $\lie{g}$, we establish non-minimality by giving a computer-assisted classification of all embeddings $\lie{h}\hookrightarrow \lie{g}$.
The proof of Theorem \ref{thm:MainTargetMinimal} is completed in this section.

\subsection{Basic properties} \label{subsect:EmbeddingBasic}

In this subsection, we establish several basic properties of embeddings $\lie{h}\hookrightarrow\lie{g}$ of type $BC'_1$ in preparation for classifying them.
Let $(\lie{g}, \lie{h})$ be a pair of type $BC'_1$ and fix $h \in \lie{h}^{-\theta}$ as in Definition \ref{def:typeBC}.
Set $\lie{a}_H \coloneq \RR h$.
Then $\lie{g}$ is simple and we have
\begin{align*}
	\lie{g} = \lie{g}_{-2} \oplus \lie{g}_{-1} \oplus \lie{g}_0 \oplus \lie{g}_1 \oplus \lie{g}_2.
\end{align*}
Note that $\lie{h}$ is isomorphic to either $\lieSU(1, k)$ ($k \geq 2$), $\lieSp(1, k)$ ($k \geq 2$) or $\lieF_{4(-20)}$.

\begin{lemma} \label{lem:NonMinimalInvariant}
	Let $(\lie{g}, \lie{h}, \orbit{O})$ be an NDD tuple.
	Assume that there exists a non-zero element $X \in \lie{g}_1 \cap \overline{\orbit{O}} \cap \lie{h}^\perp$ such that $\lie{k}_H \cap \lie{g}_0^X \neq 0$.
	Then $(\lie{g}, \lie{h}, \orbit{O})$ is not minimal.
\end{lemma}

\begin{proof}
	Take $0\neq v \in \lie{k}_H \cap \lie{g}_0^X$.
	Then $\lie{g'} \coloneq \lieCent_{\lie{g}}(v)$ is a $\theta$-stable proper subalgebra of $\lie{g}$ containing $X$ and $\lie{a}_H$.
	By definition, we have
	\begin{align*}
		\lie{g'} = \lie{h}^v \oplus (\lie{h}^\perp)^v = (\lie{g'} \cap \lie{h}) \oplus (\lie{g'} \cap \lie{h}^\perp).
	\end{align*}
	Hence $(\lie{g}, \lie{h}, \orbit{O})$ is not minimal by Proposition \ref{prop:NonMinimalSubalgebra}.
\end{proof}

We shall see that there are no minimal NDD tuples $(\lie{g}, \lie{h}, \orbit{O})$ when $\lie{h} \simeq \lieSp(1, k)$ or $\lieF_{4(-20)}$.
Assume that $\lie{h} \simeq \lieSp(1, k)$ ($k \geq 2$).
Then we have
\begin{align*}
	\lie{h}_0 \simeq \RR h \oplus \lieSp(1) \oplus \lieSp(k-1).
\end{align*}
We denote by $\lie{k}'$ the third factor $\lieSp(k-1)$.

\begin{lemma} \label{lem:SpInvariant}
	Let $V$ be an irreducible $\lie{h}$-module over $\RR$ such that the $h$-eigenspace decomposition of $V$ is written as $V = V_{-1} \oplus V_0 \oplus V_1$ with $V_1 \neq 0$.
	Then any vector in $V_1$ is $\lie{k}'$-invariant.
\end{lemma}

\begin{proof}
	It is enough to show the same assertion for an irreducible $\lie{h}$-module $V$ over $\CC$.
	Fix a Cartan subalgebra $\lieC{t}$ of $\lieSp(1+k, \CC)$ containing $h$, and take its coordinate system $\set{\varepsilon_1, \ldots, \varepsilon_{k+1}}$ such that
	\begin{align*}
		\Delta^+(\lieC{h}, \lieC{t}) &= \set{\varepsilon_i \pm \varepsilon_j : 1\leq i < j \leq k+1} \cup \set{2\varepsilon_i : 1 \leq i \leq k+1},\\
		\varepsilon_i(h) &= \begin{cases}
			1 & (i = 1,2), \\
			0 & (3 \leq i \leq k+1).
		\end{cases}
	\end{align*}
	Then a dominant integral weight in $\lieC{t}^*$ is of the form $\sum_{i=1}^{k+1} m_i \varepsilon_i$ with $m_1 \geq m_2 \geq \cdots \geq m_{k+1} \geq 0$.

	Let $\lambda$ be the highest weight of $V$.
	By assumption, we have $\lambda(h) = 1$.
	This implies that $\lambda = \varepsilon_1$; in other words, $V$ is isomorphic to the natural representation $\CC^{2k+2}$ of $\lieC{h}$.
	Hence any vector in $V_1 \simeq \CC^2$ is $\lie{k}'$-invariant.
\end{proof}

\begin{theorem} \label{thm:NonMinimalSp}
	Assume that $\lie{h} \simeq \lieSp(1, k)$ ($k \geq 2$).
	Then there exists no minimal NDD tuple $(\lie{g}, \lie{h}, \orbit{O})$.
\end{theorem}

\begin{proof}
	Let $(\lie{g}, \lie{h}, \orbit{O})$ be a minimal NDD tuple.
	By Lemma \ref{lem:ReductionLemmaStep1}, there exists a non-zero element $X \in \lie{g}_1 \cap \overline{\orbit{O}} \cap \lie{h}^\perp$.
	By Corollary \ref{cor:ReductionLemmaStep3}, we have $[X, \lie{h}_1] = 0$.
	By the grading, we have $[X, \lie{h}_2] = 0$.
	This implies that every irreducible component in the $\lie{h}$-submodule $V$ generated by $X$ satisfies the assumption of Lemma \ref{lem:SpInvariant}.
	Combining Lemmas \ref{lem:NonMinimalInvariant} and \ref{lem:SpInvariant}, we conclude that $(\lie{g}, \lie{h}, \orbit{O})$ is not minimal.
\end{proof}

Next, we assume that $\lie{h} \simeq \lieF_{4(-20)}$.
By considering the analogue of Lemma \ref{lem:SpInvariant} for this $\lie{h}$, we obtain the following stronger result.

\begin{lemma} \label{lem:NonIrreducibleF4}
	There exists no irreducible $\lie{h}$-module $V$ over $\RR$ such that the $h$-eigenspace decomposition of $V$ is written as $V = V_{-1} \oplus V_0 \oplus V_1$ with $V_1 \neq 0$.
\end{lemma}

\begin{proof}
	Fix a Cartan subalgebra $\lieC{t}$ of $\lieC{h}$ containing $h$, and take its coordinate system $\set{\varepsilon_1, \ldots, \varepsilon_4}$ such that
	\begin{align*}
		&\alpha_1 \coloneq \varepsilon_2 - \varepsilon_3, \quad \alpha_2 \coloneq \varepsilon_3 - \varepsilon_4, \\
		&\alpha_3 \coloneq \varepsilon_4, \quad \alpha_4 \coloneq \frac{1}{2}(\varepsilon_1 - \varepsilon_2 - \varepsilon_3 - \varepsilon_4)
	\end{align*}
	are simple roots in $\Delta(\lieC{h}, \lieC{t})$ and $\alpha_4(h) = 1, \alpha_i(h) = 0$ for $i = 1, 2, 3$.
	Then we have $\varepsilon_1(h) = 2$ and $\varepsilon_i(h) = 0$ for $i = 2, 3, 4$.
	The fundamental weights $\omega_1, \ldots, \omega_4$ corresponding to $\alpha_1, \ldots, \alpha_4$ are given by
	\begin{align*}
		\omega_1 &= \varepsilon_1 + \varepsilon_2, \\
		\omega_2 &= 2\varepsilon_1 + \varepsilon_2 + \varepsilon_3, \\
		\omega_3 &= \frac{1}{2}(3\varepsilon_1 + \varepsilon_2 + \varepsilon_3 + \varepsilon_4), \\
		\omega_4 &= \varepsilon_1.
	\end{align*}
	Hence any non-zero dominant integral weight $\lambda = \sum_{i=1}^4 m_i \omega_i$ with $m_i \in \ZZ_{\geq 0}$ satisfies $\lambda(h) = 2m_1 + 4m_2 + 3m_3 + 2m_4 > 1$.
	This shows the assertion.
\end{proof}

The following theorem is shown by the same argument as in the proof of Theorem \ref{thm:NonMinimalSp} using Lemma \ref{lem:NonIrreducibleF4} instead of Lemma \ref{lem:SpInvariant}.

\begin{theorem} \label{thm:NonMinimalF4}
	Assume that $\lie{h} \simeq \lieF_{4(-20)}$.
	Then there exists no minimal NDD tuple $(\lie{g}, \lie{h}, \orbit{O})$.
\end{theorem}

We shall consider the case where $\lie{h}$ is isomorphic to $\lieSU(1, k)$ ($k \geq 2$).
There are too many possible pairs $(\lie{g},\lie{h})$ to classify directly.
We give an additional condition for the pair $(\lie{g}, \lie{h})$ to have a minimal NDD tuple $(\lie{g}, \lie{h}, \orbit{O})$.

\begin{lemma} \label{lem:SUTopGradingIrr}
	$\lie{g}_2$ is an irreducible $\lie{g}_0$-module.
	In particular, $\lie{g}_2$ is a simple Jordan triple system.
\end{lemma}

\begin{proof}
	Let $\lie{g}_2 = \bigoplus_{i=1}^n V^i$ be the decomposition into irreducible $\lie{g}_0$-modules.
	Let $\lie{j}^i$ be the $\lie{g}$-submodule generated by $V^i$.
	Since the $V^i$ are annihilated by $\lie{g}_1\oplus \lie{g}_2$ and $\lie{g}_0$-stable, we have $\lie{j}^i \cap \lie{g}_2 = V^i$.
	This implies that the $\lie{j}^i$ are non-zero ideals of $\lie{g}$.
	Since $\lie{g}$ is simple, we have $n=1$ and hence $\lie{g}_2$ is irreducible.
\end{proof}

By Lemma \ref{lem:SUTopGradingIrr}, the subalgebra generated by $\lie{g}_2$ and $\theta(\lie{g}_2)$ is a simple Lie algebra with the $3$-grading induced from the grading $\lie{g} = \bigoplus_{i=-2}^2 \lie{g}_i$.
Hence the results in Subsection \ref{subsect:JordanAlgebra} can be applied to $\lie{g}_2$.

\begin{lemma} \label{lem:SUTopGrading}
	Assume that there exists a minimal NDD tuple $(\lie{g}, \lie{h}, \orbit{O})$.
	Then one has $\lie{g}_2 = \lie{h}_2$.
\end{lemma}

\begin{proof}
	By Lemma \ref{lem:ReductionLemmaStep1}, we have $\overline{\orbit{O}} \cap \lie{g}_1 \neq \set{0}$.
	Let $X$ be a non-zero element in $\overline{\orbit{O}} \cap \lie{g}_1$.
	By Lemma \ref{lem:ReductionLemmaStep3}, $\lie{g}_2^{\theta(X)}$ has codimension one in $\lie{g}_2$.
	Set $\lie{g}' \coloneq \lie{g}^X \cap \lie{g}^{\theta(X)}$.
	Then $\lie{g}'$ is a $\theta$-stable graded subalgebra of $\lie{g}$ with $\lie{g}'_2 = \lie{g}_2^{\theta(X)}$.
	
	By the definition of $\lie{g}^{\theta(X)}$, $h$ does not belong to $\lie{g'}$.
	This implies that $\lie{g}_2^{\theta(X)}$ does not contain any unitary element in $\lie{g}_2$.
	Since $\overline{\orbit{O}} \cap \lie{g}_1$ is $G_0$-stable, for any $g \in G_0$, the subspace $\Ad(g)\lie{g}_2^{\theta(X)}$ does not contain any unitary element in $\lie{g}_2$.
	By Proposition \ref{prop:CodimensionOneJordan}, we obtain $\dim(\lie{g}_2) = 1$ and hence $\lie{g}_2 = \lie{h}_2$.
\end{proof}

By Lemma \ref{lem:SUTopGrading}, we may assume that $\lie{g}_2 = \lie{h}_2$.
In other words, $(\lie{g}, \lie{h})$ is of type $BC_1$ in the sense of Definition \ref{def:typeBC}.

\begin{lemma} \label{lem:G2RootSpace}
	Let $\lie{a}$ be a maximal abelian subspace of $\lie{g}^{-\theta}$ containing $\lie{a}_H$.
	Fix a positive system $\Delta^+(\lie{g}, \lie{a})$ containing $\Delta(\lie{g}_1\oplus \lie{g}_2, \lie{a})$.
	Then $\lie{g}_2$ is the root space in $\lie{g}$ corresponding to the highest restricted root.
\end{lemma}

\begin{proof}
	Since $\dim(\lie{g}_2) = \dim(\lie{h}_2) = 1$, the $\lie{a}$-stable subspace $\lie{g}_2$ is a restricted root space with restricted root $\alpha \in \lie{a}^*$.
	By the grading, it is clear that $\alpha$ is the highest restricted root.
\end{proof}

Table \ref{table:MultiplicityHighestRoot} lists the real simple Lie algebras together with the multiplicity of the highest restricted root.
See \cite[Appendices C.3 and C.4]{Kn02}.
We omit complex Lie algebras since $\lie{g}_2$ cannot be one-dimensional.
Moreover, we exclude the case $\rank_\RR(\lie{g})=1$, as Corollary \ref{cor:SufficientSameRank} shows that no NDD tuple $(\lie{g}, \lie{h}, \orbit{O})$ can be minimal in this case.

\begin{table}[htbp]
	\centering
	\begin{tabular}{c|c}
		$\lie{g}$ & $\dim(\lie{g}_{2})$ \\
		\hline
		$\lieSL(n, \mathbb{R})$ ($n \geq 3$) & $1$ \\
		$\lieSU(p, q)$ ($2 \leq p \leq q$) & $1$ \\
		$\lieSL(n, \HH)$ ($n \geq 2$) & $4$ \\
		$\lieSO(p, q)$ ($2 \leq p \leq q$) & $1$ \\
		$\lieSO^*(2n)$ ($n \geq 4$) & $1$ \\
		$\lieSp(n, \mathbb{R})$ ($n \geq 3$) & $1$ \\
		$\lieSp(p, q)$ ($2 \leq p \leq q$) & $3$ \\
		$\lieE_{6(6)}$ & $1$ \\
		$\lieE_{6(2)}$ & $1$ \\
		$\lieE_{6(-14)}$ & $1$ \\
		$\lieE_{6(-26)}$ & $8$ \\
		$\lieE_{7(7)}$ & $1$ \\
		$\lieE_{7(-5)}$ & $1$ \\
		$\lieE_{7(-25)}$ & $1$ \\
		$\lieE_{8(8)}$ & $1$ \\
		$\lieE_{8(-24)}$ & $1$ \\
		$\lieG_{2(2)}$ & $1$ \\
		$\lieF_{4(4)}$ & $1$ \\
	\end{tabular}
	\caption{Real simple Lie algebras $\lie{g}$ and the multiplicity of the highest restricted root}
	\label{table:MultiplicityHighestRoot}
\end{table}

To describe the embedding of $\lieSU(1, k)$, we examine some properties of the fundamental representations of $\lieSU(1, k)$.

\begin{lemma} \label{lem:FundamentalRep}
	Let $V$ be a non-trivial irreducible representation of $\lie{h} = \lieSU(1, k)$ over a field $\KK \in \set{\RR, \CC, \HH}$.
	Assume that the eigenspace decomposition of $V$ with respect to the action of $h$ is written as $V = V_{-1} \oplus V_0 \oplus V_1$, where $V_i$ is the eigenspace of $h$ with eigenvalue $i$.
	\begin{enumparen}
		\item If $\KK = \CC$, then $V$ is isomorphic to $\bigwedge\nolimits^j \CC^{k+1}$ for some $0 < j \leq k$.
		\item If $\KK = \CC$ and $\dim_{\CC}(V_1) = 1$, then one has $j = 1 \text{or} k$ in (1).
		If $\KK = \CC$ and $\dim_{\CC}(V_1) = 2$, then one has $(j, k) = (2, 3)$.
		\item $\KK = \RR$ and $\dim_{\RR}(V_1) = 1$ cannot happen.
		\item If $\KK = \HH$ and $\dim_{\HH}(V_1) = 1$, then there exist two possible cases:
		\begin{align*}
			V&\simeq \HH\otimes_{\CC}\CC^{k+1}, \\
			V&\simeq \HH^3 \quad \text{under the isomorphism $\lieSU(1, 3) \simeq \lieSO^*(6)$},
		\end{align*}
		where the second case occurs when $k = 3$.
		\item $\bigwedge\nolimits^2 \CC^4$ has no real form.
	\end{enumparen}
\end{lemma}

\begin{proof}
	Realize $\lie{h}_{\CC}$ as the standard matrix Lie algebra $\lie{sl}(k+1, \CC)$ such that $h$ is represented by
	\begin{align*}
		\begin{pmatrix}
			1 & 0 & \cdots & 0 & 0 \\
			0 & 0 & \cdots & 0 & 0 \\
			\vdots & \vdots & \ddots & \vdots & \vdots \\
			0 & 0 & \cdots & 0 & 0 \\
			0 & 0 & \cdots & 0 & -1
		\end{pmatrix}.
	\end{align*}
	Consider the standard Borel subalgebra and Cartan subalgebra of $\lie{sl}(k+1, \CC)$ consisting of upper triangular matrices and diagonal matrices, respectively.
	Then the highest weight of $V$ is a fundamental weight $\omega_j$ ($1 \leq j \leq k$) by the assumption $V = V_1 \oplus V_0 \oplus V_{-1}$.
	In other words, for each $V$, there exists $0 < j \leq k$ such that $V \simeq \bigwedge\nolimits^j \CC^{k+1}$.
	This shows (1).

	For each $V = \bigwedge\nolimits^j \CC^{k+1}$ with $0 < j \leq k$, we have $V_1 \simeq \bigwedge\nolimits^{j-1} \CC^{k-1}$.
	Hence $\dim_{\CC}(V_1) \in \set{1, 2}$ if and only if $j \in \set{1, k}$, or $(j, k) = (2, 3)$.
	This shows (2).

	To show (3), assume that $\KK = \RR$ and $\dim_{\RR}(V_1) = 1$.
	Then $V\otimes_{\RR} \CC$ is irreducible over $\CC$ and $\dim_{\CC}(V_1) = 1$.
	By (2), we have $V\otimes_{\RR} \CC \simeq \CC^{k+1}$ or $(\CC^{k+1})^*$.
	By definition, $V\otimes_{\RR}\CC$ has a real form.
	This happens only if $\lieSU(1, k) \simeq \lieSL(1+k, \RR)$.
	This contradicts the assumption $k \geq 2$.

	We shall show (4).
	Take an irreducible subrepresentation $W$ of $V$ over $\RR$.
	Since $\HH\cdot W = V$, any irreducible subrepresentation of $V$ over $\RR$ is isomorphic to $W$.
	Hence $V$ is completely reducible over $\RR$.
	By (3), we have $\dim_{\RR}(W \cap V_1) = 2$ or $4$.
	
	Since $\HH\cdot W = V$, we have
	\begin{align*}
		V \simeq \Hom_{\RR, \lie{h}}(W, V) \otimes_{\End_{\RR, \lie{h}}(W)} W.
	\end{align*}
	Since $V$ is irreducible over $\HH$, the space $\Hom_{\RR, \lie{h}}(W, V)$ is one-dimensional over $\HH$.
	Comparing the dimensions, we have
	\begin{align*}
		\End_{\RR, \lie{h}}(W) \simeq \begin{cases}
			\CC & (\dim_{\RR}(W \cap V_1) = 2), \\
			\HH & (\dim_{\RR}(W \cap V_1) = 4).
		\end{cases}
	\end{align*}
	This implies that $W$ has a complex structure.
	By (2), we obtain
	\begin{align*}
		W \simeq \begin{cases}
			\CC^{k+1} & (\dim_{\RR}(W \cap V_1) = 2), \\
			\bigwedge\nolimits^2 \CC^4 & (\dim_{\RR}(W \cap V_1) = 4).
		\end{cases}
	\end{align*}
	Note that $\CC^{k+1}$ is isomorphic to $(\CC^{k+1})^* \simeq \overline{\CC^{k+1}}$ as a representation over $\RR$.
	Therefore we have shown (4).

	(5) follows from the proof of (4) above.
	In fact, if $\bigwedge\nolimits^2 \CC^4$ has a real form, then the real form is isomorphic to $\CC^4$ by the proof of (4), which is a contradiction.
	Note that any finite-dimensional $\HH$-representation of a simple real Lie algebra that is irreducible over $\CC$ cannot admit a real form.
\end{proof}

Classifying embeddings of $\lie{h}$ is not easy if one only uses the fact that they are embeddings of $\lieSU(1,k)$.
However, by using the fact that $\lieC{h}$ is a regular subalgebra of $\lieC{g}$, as shown in the following lemma, the classification can be reduced to computations in root systems.

\begin{lemma} \label{lem:EmbeddingSU}
	Assume that $\lie{h} \simeq \lieSU(1, k)$ ($k \geq 2$).
	Then $\lie{g}$ has a compact Cartan subalgebra contained in $\lieNorm_{\lie{g}}(\lie{h})$.
	In particular, $\lieC{h}$ is a regular subalgebra of $\lieC{g}$.
\end{lemma}

\begin{proof}
	Decompose the $\lieC{h}$-module $\lieC{g}$ into irreducible components as
	\begin{align*}
		\lieC{g} = \lieC{h} \oplus \lieCent_{\lieC{g}}(\lieC{h}) \oplus \bigoplus_{i=1}^m V_i,
	\end{align*}
	where $V_i$ are irreducible $\lieC{h}$-modules with non-trivial $\lieC{h}$-action.
	Since $\lie{g}_2 = \lie{h}_2$, we have 
	\begin{align*}
		\bigoplus_{i=1}^m V_i \subset (\lieC{g})_{-1} \oplus (\lieC{g})_0 \oplus (\lieC{g})_1.
	\end{align*}

	Fix a compact Cartan subalgebra $\lie{t}$ of $\lie{h}$.
	By Lemma \ref{lem:FundamentalRep}, for each $V_i$, there exists $0 < j \leq k$ such that $V_i \simeq \bigwedge\nolimits^j \CC^{k+1}$.
	In particular, none of the $V_i$ has zero weight, and we have $V_i^{\lie{t}} = 0$.
	This implies that
	\begin{align}
		(\lieC{g})^{\lie{t}} = (\lieC{h})^{\lie{t}} \oplus \lieCent_{\lieC{g}}(\lieC{h}) \oplus \bigoplus_{i=1}^m V_i^{\lie{t}} = \lieC{t} \oplus \lieCent_{\lieC{g}}(\lieC{h}). \label{eqn:CentralizerCartan}
	\end{align}
	Since $(\lie{g}, \lie{h})$ is of type $BC_1$, $\lieCent_{\lie{g}}(\lie{h})$ is compact.
	Take a maximal torus $\lie{t'}$ of $\lieCent_{\lie{g}}(\lie{h})$.
	Then $\lie{t}\oplus \lie{t'}$ is a compact Cartan subalgebra of $\lie{g}$ by \eqref{eqn:CentralizerCartan}.
\end{proof}

Combining Lemma \ref{lem:EmbeddingSU} and Table \ref{table:MultiplicityHighestRoot}, we obtain the following corollary.

\begin{corollary} \label{cor:EmbeddingSU}
	If $\lie{h} \simeq \lieSU(1, k)$ ($k \geq 2$), then $\lie{g}$ is isomorphic to one of the following Lie algebras:
	\begin{align*}
		&\lieSU(p, q) \quad (p, q \geq 1, p+q\geq 3), \\
		&\lieSO(p, q) \quad (p, q \geq 2, p+q\geq 7, p\text{:even}), \\
		&\lieSO^*(2n) \quad (n \geq 5), \\
		&\lieSp(n,\RR) \quad (n \geq 2), \\
		&\lieE_{6(2)}, \quad \lieE_{6(-14)}, \\
		&\lieE_{7(7)}, \quad \lieE_{7(-5)}, \quad \lieE_{7(-25)}, \\
		&\lieE_{8(8)}, \quad \lieE_{8(-24)}, \\
		&\lieF_{4(4)}, \\
		&\lieG_{2(2)}.
	\end{align*}
\end{corollary}

\subsection{Embeddings of \texorpdfstring{$\lieSU(1, k)$}{su(1,k)}: Classical \texorpdfstring{$\lie{g}$}{g}}

In this subsection, we concentrate on the case where $\lie{h}$ is isomorphic to $\lieSU(1, k)$ ($k \geq 2$) and $\lie{g}$ is classical.
By Corollary \ref{cor:EmbeddingSU}, the possible choices of $\lie{g}$ are listed as follows:
\begin{align*}
	&\lieSU(p, q) \quad (p, q \geq 2), \\
	&\lieSO(p, q) \quad (p, q \geq 2, p+q\geq 7, p\text{:even}), \\
	&\lieSO^*(2n) \quad (n \geq 5), \\
	&\lieSp(n,\RR) \quad (n \geq 2).
\end{align*}
For simplicity, we consider $\lieU(p, q)$ instead of $\lieSU(p, q)$.
Although $(\lie{g}, \lie{h})$ is not of type $BC_1$ if $\lie{g} = \lieU(p, q)$, this replacement does not affect the classification of embeddings $\lie{h}\hookrightarrow \lie{g}$.
If $\lie{g}$ has real rank one, then any NDD tuple $(\lie{g}, \lie{h}, \orbit{O})$ is not minimal by Corollary \ref{cor:SufficientSameRank}, so we omit the case.

Let $V\coloneq \KK^n$ be the natural representation of $\lie{g}$, where $\KK = \RR$, $\CC$ or $\HH$.
Then $\lie{g}$ is isomorphic to the Lie algebra of the isometry group of $V$ with respect to some non-degenerate bilinear or sesquilinear form $B(\cdot, \cdot)$.
Since $h$ is the coroot of the highest restricted root, $V$ has the eigenspace decomposition
\begin{align}
	V &= V_{-1} \oplus V_0 \oplus V_1, \\
	\dim_{\KK}(V_1) &= \begin{cases}
		2 & (\lie{g} = \lieSO(p, q)) \\
		1 & (\text{otherwise})
	\end{cases} \label{eqn:eigenspacehSU}
\end{align}
with respect to the action of $h$.
Then the restriction of $B(\cdot, \cdot)$ to $V_1\oplus V_{-1}$ is non-degenerate.

Set $W \coloneq (V^{\lie{h}})^\perp$.
Then $V = W \oplus V^{\lie{h}}$.
By \eqref{eqn:eigenspacehSU} and Lemma \ref{lem:FundamentalRep}, $W$ is irreducible over $\KK$.

\begin{lemma} \label{lem:EmbeddingSUtoSp}
	$\lie{g}$ is not isomorphic to $\lieSp(n, \RR)$.
\end{lemma}

\begin{proof}
	By Lemma \ref{lem:FundamentalRep}, $\KK = \RR$ and $\dim_{\RR}(V_1) = 1$ cannot occur.
	This implies that $\lie{g}$ is not isomorphic to $\lieSp(n, \RR)$.
\end{proof}

\begin{lemma} \label{lem:EmbeddingSUtoSU}
	Suppose that $\lie{g} = \lieU(p, q)$.
	Then one has $\lieNorm_{\lie{g}}(\lie{h}) = \lieU(1, q) \oplus \lieU(p-1, 0)$ or $\lieU(p, 1) \oplus \lieU(0, q-1)$ up to automorphisms of $\lie{g}$.
	In particular, $\lieNorm_{\lie{g}}(\lie{h})$ is a symmetric subalgebra of $\lie{g}$.
\end{lemma}

\begin{proof}
	We have
	\begin{align*}
		\lie{h} \subset \lieU(W) \oplus \lieU(V^{\lie{h}}) \subset \lieU(V) = \lie{g}.
	\end{align*}
	Since $\lieCent_{\lie{g}}(\lie{h})$ is compact, $B$ is definite on $V^{\lie{h}}$.
	By Lemma \ref{lem:FundamentalRep}, $W$ is isomorphic to $\CC^{k+1}$ or $(\CC^{k+1})^*$ as an $\lie{h}$-module over $\CC$.
	Hence we have $\lie{h} = \lieSU(W)$.
	Therefore, we obtain $\lieNorm_{\lie{g}}(\lie{h}) = \lieU(W) \oplus \lieU(V^{\lie{h}})$.
\end{proof}

\begin{lemma} \label{lem:EmbeddingSUtoSOstar}
	Suppose that $\lie{g} = \lieSO^*(2n)$.
	Then one has
	\begin{align*}
		\lieNorm_{\lie{g}}(\lie{h}) &= \lieU(1, n-1) \quad (n = k+1), \text{ or }\\
		\lieNorm_{\lie{g}}(\lie{h}) &= \lieU(1, n-2) \oplus \lieSO^*(2) \quad (n = k+2),
	\end{align*}
	up to automorphisms of $\lie{g}$.
\end{lemma}

\begin{proof}
	We have
	\begin{align*}
		\lie{h} \subset \lieSO^*(W) \oplus \lieSO^*(V^{\lie{h}}) \subset \lieSO^*(V) = \lie{g}.
	\end{align*}
	Since $\lieCent_{\lie{g}}(\lie{h})$ is compact, we have $\dim_{\HH}(V^{\lie{h}}) \leq 1$.
	By Lemma \ref{lem:FundamentalRep}, $W$ is isomorphic to one of the following representations:
	\begin{align*}
		\HH\otimes_{\CC} \CC^{k+1}, \quad \HH^3 \left(\simeq \bigwedge\nolimits^2 \CC^4 \right).
	\end{align*}
	If $W \simeq \HH^3$, then we have $n \leq 4$, and this contradicts the assumption $n \geq 5$.

	Assume that $W \simeq \HH\otimes_{\CC} \CC^{k+1}$.
	Then the embedding is
	\begin{align*}
		\lieSU(1, k) \simeq \lie{h} \hookrightarrow \lieSO^*(W) \hookrightarrow \lieSO^*(V).
	\end{align*}
	If $V^{\lie{h}} = 0$, then we have $n=k+1$ and $\lieNorm_{\lie{g}}(\lie{h}) \simeq \lieU(1, n-1)$, which is a symmetric subalgebra of $\lie{g}$.
	If $\dim_{\HH}(V^{\lie{h}}) = 1$, then we have $n = k+2$ and $\lieNorm_{\lie{g}}(\lie{h}) \simeq \lieU(1, n-2) \oplus \lieSO^*(2)$, which is a symmetric subalgebra of $\lieSO^*(W) \oplus \lieSO^*(V^{\lie{h}})\subset \lie{g}$.
	This shows the assertion.
\end{proof}

\begin{lemma} \label{lem:EmbeddingSUtoSO}
	Suppose that $\lie{g} = \lieSO(p, q)$.
	Then, up to automorphisms of $\lie{g}$, one has
	\begin{align*}
		\lieNorm_{\lie{g}}(\lie{h}) &= \lieU(p/2, 1) \oplus \lieSO(0, q - 2),  \\
		\lieNorm_{\lie{g}}(\lie{h}) &= \lieU(1, q/2) \oplus \lieSO(p - 2, 0) \quad (q\text{: even}), \text{ or } \\
		\lieNorm_{\lie{g}}(\lie{h}) &= \lieU(1, k) \oplus \lieSO(0, q - 2k) \quad (p=2, 2k \leq q).
	\end{align*}
\end{lemma}

\begin{proof}
	We have
	\begin{align*}
		\lie{h} \subset \lieSO(W) \oplus \lieSO(V^{\lie{h}}) \subset \lieSO(V) = \lie{g}.
	\end{align*}
	Since $\lieCent_{\lie{g}}(\lie{h})$ is compact, $B$ is definite on $V^{\lie{h}}$.
	Consider the representation $W\otimes_{\RR} \CC$ over $\CC$.
	Since $\dim_{\RR}(V_1) = 2$, Lemma \ref{lem:FundamentalRep} (2) implies that $W\otimes_{\RR} \CC$ is reducible or isomorphic to $\bigwedge\nolimits^2 \CC^{4}$.
	By Lemma \ref{lem:FundamentalRep} (5), $\bigwedge\nolimits^2 \CC^{4}$ has no real form, and hence $W\otimes_{\RR} \CC$ is reducible.
	In other words, $W$ has a complex structure.
	By Lemma \ref{lem:FundamentalRep} (2), $W$ is isomorphic to $\CC^{k+1}$ or $(\CC^{k+1})^*$.

	Since $((W \otimes_\RR W)^*)^{\lie{h}}$ is spanned by the symmetric bilinear form $B|_W$ and some skew-symmetric bilinear form, the complex structure on $W$ is compatible with $B|_W$.
	Hence we have $\lie{h} = \lieSU(W)$ and $(\lieSO(W), \lieU(W))$ is a symmetric pair.
	Therefore, we obtain $\lieNorm_{\lie{g}}(\lie{h}) = \lieU(W) \oplus \lieSO(V^{\lie{h}})$.
\end{proof}

We have thus proved non-minimality in the cases where $\lie{g}=\lieSp(n,\RR)$ or $\lie{g}=\lieSU(p,q)$.
The cases where $\lie{g}=\lieSO(p,q)$ or $\lie{g}=\lieSO^*(2n)$ will be treated in the next two subsections.

\subsection{Explicit computation for \texorpdfstring{$\lie{g} = \lieSO(p,q)$}{g=SO(p,q)}}

 We shall show that $(\lie{g}, \lie{h}, \orbit{O})$ is not minimal using Lemma \ref{lem:ReductionRealStronglyOrthogonal} when $\lie{g}$ is isomorphic to $\lieSO(p,q)$.
We check the assumptions of Lemma \ref{lem:ReductionRealStronglyOrthogonal} directly.
By Lemma \ref{lem:EmbeddingSUtoSO}, it is enough to consider the following cases.
\begin{enumerate}[label=(\alph*)]
	\item $(\lie{g}, \lieNorm_{\lie{g}}(\lie{h})) = (\lieSO(p, q), \lieU(p/2, 1)\oplus \lieSO(0, q-2))$ ($p\geq 4, q \geq 4$, $p$: even)
	\item $(\lie{g}, \lieNorm_{\lie{g}}(\lie{h})) = (\lieSO(p, 3), \lieU(p/2, 1))$ ($p\geq 4$, $p$: even)
	\item $(\lie{g}, \lieNorm_{\lie{g}}(\lie{h})) = (\lieSO(2, q), \lieU(1, k)\oplus \lieSO(0, q-2k))$ ($q > 2k \geq 4$)
\end{enumerate}
Note that, if $q = 2k$ in case (c), then $(\lie{g}, \lieNorm_{\lie{g}}(\lie{h}))$ is a symmetric pair and hence $(\lie{g}, \lie{h}, \orbit{O})$ is not minimal for any $\orbit{O}$ by Lemma \ref{lem:SufficientSymmetric}.

Let $V \coloneq \RR^{p+q}$ with the standard inner product of signature $(p, q)$ ($p\geq 4, q \geq 3$, $p$: even).
Let $\set{e_1, \ldots, e_{p+q}}$ be the standard basis of $V$.
Set
\begin{align*}
	V_1 &\coloneq \spn{\RR}{e_1 + e_{p+q}, e_2 - e_{p+q-1}}, \\
	V_{-1} &\coloneq \spn{\RR}{e_1 - e_{p+q}, e_2 + e_{p+q-1}}, \\
	W_0 &\coloneq \spn{\RR}{e_3, \ldots, e_{p}}, \\
	W &\coloneq V_{-1} \oplus W_0 \oplus V_{1}, \\
	W' &\coloneq \spn{\RR}{e_{p+1}, \ldots, e_{p+q-2}}, \\
	V_0 &\coloneq W_0 \oplus W'.
\end{align*}
Then $V_1$ and $V_{-1}$ are isotropic subspaces of $V$ and we have $V = W \oplus W'$.
Define a complex structure on $W$ by
\begin{align*}
	J(e_{2i-1}) &= e_{2i} &&(1\leq i \leq p/2), \\
	J(e_{2i}) &= -e_{2i-1} &&(1\leq i \leq p/2), \\
	J(e_{p+q-1}) &= e_{p+q}, \\
	J(e_{p+q}) &= -e_{p+q-1}.
\end{align*}
$J$ is skew-symmetric on $W$ and hence $W$ is a complex vector space with a Hermitian form of signature $(p/2, 1)$.
Set $\lie{g} = \lieSO(V)$ and $\lie{h} = \lieSU(W)$.

In this setting, we may choose $h \in \lie{h}$ such that $h$ acts on $V_{\pm 1}$ by $\pm 1$ and on $V_0$ by $0$.
The Levi subalgebra $\lie{l} \coloneq \lieCent_{\lie{g}}(h)$ is $\lieGL(V_1) \oplus \lieSO(V_0)$.
We identify the $\lie{l}$-module $\lie{g}_1$ with $\Hom_{\RR}(V_{-1}, V_0)$.
Then we have
\begin{align*}
	\lie{g}_1 \cap \lie{h} &= \Hom_{\CC}(V_{-1}, W_0), \\
	\lie{g}_1 \cap \lie{h}^\perp &= \Hom_{\CC}(V_{-1}, \overline{W_0}) \oplus \Hom_{\RR}(V_{-1}, W').
\end{align*}

First, we consider the case of $q \geq 4$.
Define $T_1, \ldots, T_4 \in \Hom_{\RR}(V_{-1}, V_0)$ by
\begin{align*}
	T_1(e_1 - e_{p+q}) &= e_3 + e_{p+1}, & T_1(e_2 + e_{p+q-1}) &= 0, \\
	T_2(e_1 - e_{p+q}) &= e_{3} - e_{p+1}, & T_2(e_2 + e_{p+q-1}) &= 0, \\
	T_3(e_1 - e_{p+q}) &= 0, & T_3(e_2 + e_{p+q-1}) &= -e_4 - e_{p+2}, \\
	T_4(e_1 - e_{p+q}) &= 0, & T_4(e_2 + e_{p+q-1}) &= -e_4 + e_{p+2},
\end{align*}
and set $\lie{c} \coloneq \spn{\RR}{T_1, T_2, T_3, T_4}$.

\begin{lemma} \label{lem:SOpq4Basic}
	Let $a, b, c, d \in \RR$.
	\begin{enumparen}
		\item There exists a maximal abelian subspace $\lie{a}$ of $\lie{g}^{-\theta}$ containing $h$ such that $T_1, \ldots, T_4$ are root vectors with respect to $\lie{a}$ and the roots are strongly orthogonal.
		\item $aT_1 + bT_2 + cT_3 + dT_4 \in \lie{h}^\perp$ if and only if $a + b = c + d$.
		\item $aT_1 + bT_2 + cT_3 + dT_4 \in \lie{h}$ if and only if $a = b$, $c = d$ and $a + c = 0$.
		\item $\lie{c} = (\lie{c} \cap \lie{h}) \oplus (\lie{c} \cap \lie{h}^\perp)$.
	\end{enumparen}
\end{lemma}

\begin{proof}
	The assertions are straightforward from the definitions of $T_1, \ldots, T_4$.
\end{proof}

We show that there is no minimal tuple by applying Lemma \ref{lem:ReductionRealStronglyOrthogonal}.
The condition $\lie{c}=(\lie{c} \cap \lie{h}) \oplus (\lie{c} \cap \lie{h}^\perp)$ has already been verified in Lemma \ref{lem:SOpq4Basic}.
It remains to verify the other condition, namely $\overline{\orbit{O}} \cap \lie{h}^\perp \cap \lie{c} \neq \set{0}$.
Note that $\lie{c}$ is contained in $\Nilpotent(\lie{g}, H)$ since $\lie{c} \subset \lie{g}_1$.

\begin{lemma} \label{lem:SOpq4}
	Let $\orbit{O}$ be a nilpotent orbit in $\lie{g}$.
	Assume that there exists $T \in \lie{g}_1 \cap \lie{h}^\perp \cap \overline{\orbit{O}}$ such that the component of $T$ in $\Hom_{\CC}(V_{-1}, \overline{W_0})$ is non-zero.
	Then one of $T_1 + T_3$ and $T_1 + T_4$ belongs to $\overline{\orbit{O}} \cap \lie{h}^\perp$.
\end{lemma}

\begin{proof}
	By assumption, we have
	\begin{align*}
		T(e_1 - e_{p+q}) &= v + w, \\
		T(e_2 + e_{p+q-1}) &= -J(v) + w',
	\end{align*}
	for some $0\neq v \in W_0$, $w \in W'$, and $w' \in W'$.
	Recall that $\overline{\orbit{O}}$ is a $G$-stable closed cone.
	By the action of $\LieU(W_0)$ and scalar multiplication, we may assume that $v = e_3$.
	Then we have $T(e_2 + e_{p+q-1}) = -J(e_3) + w' = -e_4 + w'$.
	By the action of $\LieSO(W')$, we may assume that $w = ae_{p+1}$ and $w' = be_{p+1} + ce_{p+2}$ for some $a \in \RR_{\geq 0}$, $b, c \in \RR$.
	
	Define an element $h' \in \lieGL(V_1) \oplus \lieSO(V_0)$ by
	\begin{align*}
		h' \cdot (e_1 - e_{p+q}) &= e_1 - e_{p+q}, \\
		h'\cdot (e_2 + e_{p+q-1}) &= 2(e_2 + e_{p+q-1}), \\
		h'\cdot (e_3 + e_{p+1}) &= e_3 + e_{p+1}, \\
		h'\cdot (e_3 - e_{p+1}) &= -(e_3 - e_{p+1}), \\
		h'\cdot (-e_4 + e_{p+2}) &= 2\varepsilon (-e_4 + e_{p+2}), \\
		h'\cdot (-e_4 - e_{p+2}) &= -2\varepsilon(-e_4 - e_{p+2}), \\
		h'\cdot v &= 0 \quad (v \in (V_{-1} \oplus V_1 \oplus \spn{\RR}{e_3, e_4, e_{p+1}, e_{p+2}})^\perp),
	\end{align*}
	where $\varepsilon = 1$ if $c \geq 0$ and $\varepsilon = -1$ if $c < 0$.
	Set $T'\coloneq \lim_{t\to \infty} (\exp(th')T)$.
	Then we have
	\begin{align*}
		T'(e_1 - e_{p+q}) &= \lim_{t\to \infty} e^{-t}\exp(th')(e_3 + ae_{p+1})\\ 
		&= \frac{1+a}{2}(e_3 + e_{p+1}), \\
		T'(e_2 + e_{p+q-1}) &= \lim_{t\to \infty} e^{-2t}\exp(th')(-e_4 + be_{p+1} + ce_{p+2}) \\
		&= \begin{cases}
			\frac{1+c}{2}(-e_4 + e_{p+2}) & (c \geq 0), \\
			\frac{1-c}{2}(-e_4 - e_{p+2}) & (c < 0).
		\end{cases}
	\end{align*}
	Hence one of $T_1 + T_3$ and $T_1 + T_4$ belongs to $\LieGL(V_1)_o T' \subset \overline{\orbit{O}}$.
	This shows the assertion.
\end{proof}

\begin{theorem} \label{thm:SOpq4}
	If $\lie{g} \simeq \lieSO(p, q)$ ($p\geq 4, q \geq 4$, $p$: even) and $\lie{h} \simeq \lieSU(p/2, 1)$, then $(\lie{g}, \lie{h}, \orbit{O})$ is not minimal for any $\orbit{O}$.
\end{theorem}

\begin{proof}
	To apply Lemma \ref{lem:ReductionRealStronglyOrthogonal}, we shall show that $\lie{c} \cap \lie{h}^\perp \cap \overline{\orbit{O}} \neq \set{0}$.
	If the assumption of Lemma \ref{lem:SOpq4} holds, then we have nothing to prove.
	Assume that there exists $0\neq T \in \lie{g}_1 \cap \lie{h}^\perp \cap \overline{\orbit{O}}$ such that $T\in \Hom_{\RR}(V_{-1}, W')$.

	Recall that $\Hom_{\RR}(V_{-1}, W')$ is a $\LieGL(V_1)_o\times \LieSO(W')$-stable subspace in $\lie{g}_1\cap \lie{h}^\perp$ and the bilinear form on $W'$ is definite.
	Hence there exists $g \in \LieGL(V_1)_o\times \LieSO(W')$ such that $(gT)(e_1 - e_{p+q}) = ae_{p+1}$ and $(gT)(e_2 + e_{p+q-1}) = be_{p+2}$ for some $a, b \in \RR$.
	Then we have $0\neq gT \in \lie{c} \cap \lie{h}^\perp \cap \overline{\orbit{O}}$.
	We have shown the theorem.
\end{proof}

Next, we assume that $q = 3$.
In this case, Lemma \ref{lem:SOpq4} fails since $\dim(W') = 1$.
Instead, we prove an analogous lemma and apply Lemma \ref{lem:ReductionRealStronglyOrthogonal}.
Since the argument follows almost the same outline, we omit the details.

Define $T_1, \ldots, T_3 \in \Hom_{\RR}(V_{-1}, V_0)$ by
\begin{align*}
	T_1(e_1 - e_{p+q}) &= e_3 + e_{p+1}, & T_1(e_2 + e_{p+q-1}) &= 0, \\
	T_2(e_1 - e_{p+q}) &= e_{3} - e_{p+1}, & T_2(e_2 + e_{p+q-1}) &= 0, \\
	T_3(e_1 - e_{p+q}) &= 0, & T_3(e_2 + e_{p+q-1}) &= -e_4,
\end{align*}
and set $\lie{c} \coloneq \spn{\RR}{T_1, T_2, T_3}$.
Then $\lie{c}$ satisfies properties analogous to those in Lemma \ref{lem:SOpq4Basic}.
In particular, there exists a maximal abelian subspace $\lie{a}$ of $\lie{g}^{-\theta}$ containing $h$ such that
$T_1, \ldots, T_3$ are root vectors with respect to $\lie{a}$ and the roots are strongly orthogonal.

For $0 \neq T \in \lie{g}_1 \cap \lie{h}^\perp$, there are four possibilities:
\begin{enumparen}
	\item $B$ is positive definite on $\Im(T)$.
	\item $B$ is negative definite on $\Im(T)$ and $\dim(\Im(T)) = 1$.
	\item $B$ has signature $(1, 1)$ on $\Im(T)$.
	\item $B$ is degenerate and positive semi-definite on $\Im(T)$.
\end{enumparen}
We have $\dim(\Im(T)) = 2$ if and only if the component of $T$ in $\Hom_{\CC}(V_{-1}, \overline{W_0})$ is non-zero.

\begin{lemma} \label{lem:SOpq3-degenerate}
	Let $\orbit{O}$ be a nilpotent orbit in $\lie{g}$.
	Assume that there exists $0 \neq T \in \lie{g}_1 \cap \lie{h}^\perp \cap \orbit{O}$ such that $\dim(\Im(T)) = 2$.
	Then one has $\lie{c} \cap \lie{h}^\perp \cap \orbit{O} \neq \set{0}$.
\end{lemma}

\begin{proof}
	Since the projection from $\Im(T)$ to $W'$ has a non-trivial kernel, there exists $0\neq v \in V_{-1}$ such that $T(v) \in W_0$.
	By the action of $\LieGL(V_1, J)\times \LieU(W_0) \subset \LieU(W)$, we may assume that $T(e_2 + e_{p+q-1}) = -e_4$, where $\LieGL(V_1, J)$ is the general linear group of the complex vector space $(V_1, J)$.
	Then we have $T(e_1 - e_{p+q}) = -J(e_4) + ae_{p+1} = e_3 + ae_{p+1}$ for some $a \in \RR$.
	This shows the assertion.
\end{proof}

\begin{theorem} \label{thm:SOpq3}
	If $\lie{g} \simeq \lieSO(p, 3)$ and $\lie{h} \simeq \lieSU(p/2, 1)$, then $(\lie{g}, \lie{h}, \orbit{O})$ is not minimal for any $\orbit{O}$.
\end{theorem}

\begin{proof}
	The proof is the same as the proof of Theorem \ref{thm:SOpq4} except that we use Lemma \ref{lem:SOpq3-degenerate} instead of Lemma \ref{lem:SOpq4}.
\end{proof}

We shall consider the remaining case where $\lie{g} \simeq \lieSO(2, q)$ ($q \geq 5$).
Let $V \coloneq \RR^{2+q}$ with the standard inner product of signature $(2, q)$ ($q \geq 5$).
Let $2\leq k < q/2$.
Let $\set{e_1, \ldots, e_{2+q}}$ be the standard basis of $V$.
Set
\begin{align*}
	V_1 &\coloneq \spn{\RR}{e_1 + e_{2+q}, e_2 - e_{2+q-1}}, \\
	V_{-1} &\coloneq \spn{\RR}{e_1 - e_{2+q}, e_2 + e_{2+q-1}}, \\
	W_0 &\coloneq \spn{\RR}{e_3, \ldots, e_{2k}}, \\
	W &\coloneq V_{-1} \oplus W_0 \oplus V_{1}, \\
	W' &\coloneq \spn{\RR}{e_{2k+1}, \ldots, e_{2+q-2}}, \\
	V_0 &\coloneq W_0 \oplus W'.
\end{align*}
Then $V_1$ and $V_{-1}$ are isotropic subspaces of $V$ and we have $V = W \oplus W'$.
Define a complex structure on $W$ by
\begin{align*}
	J(e_{2i-1}) &= e_{2i} &&(1\leq i \leq k), \\
	J(e_{2i}) &= -e_{2i-1} &&(1\leq i \leq k), \\
	J(e_{2+q-1}) &= e_{2+q}, \\
	J(e_{2+q}) &= -e_{2+q-1}.
\end{align*}
$J$ is skew-symmetric on $W$ and hence $W$ is a complex vector space with a Hermitian form of signature $(1, k)$.
Set $\lie{g} = \lieSO(V)$ and $\lie{h} = \lieSU(W)$.

Take $h \in \lie{h}$ and identify the $\lie{l}$-module $\lie{g}_1$ with $\Hom_{\RR}(V_{-1}, V_0)$ as in the previous cases.
Define $T_1, \ldots, T_3 \in \Hom_{\RR}(V_{-1}, V_0)$ by
\begin{align*}
	T_1(e_1 - e_{2+q}) &= e_3, & T_1(e_2 + e_{2+q-1}) &= 0, \\
	T_2(e_1 - e_{2+q}) &= 0, & T_2(e_2 + e_{2+q-1}) &= e_4, \\
	T_3(e_1 - e_{2+q}) &= e_{2k+1}, & T_3(e_2 + e_{2+q-1}) &= 0
\end{align*}
and set $\lie{c} \coloneq \spn{\RR}{T_1, T_2}$.
Then $\lie{c}$ satisfies properties analogous to those in Lemma \ref{lem:SOpq4Basic}.
In particular, there exists a maximal abelian subspace $\lie{a}$ of $\lie{g}^{-\theta}$ containing $h$ such that
$T_1, \ldots, T_3$ are root vectors with respect to $\lie{a}$ and $\set{T_1, T_2}$ is a $\lie{q}$-orthogonal system defined in \ref{subsect:ReductionStronglyOrthogonal}.

\begin{lemma} \label{lem:SOpq2-nondegenerate}
	Let $T \in \lie{g}_1$ such that $\dim(\Im(T)) = 2$.
	Then there exists $g \in G$ such that $\Ad(g)T \in \lie{c} \cap \lie{h}^\perp$.
\end{lemma}

\begin{proof}
	Since the bilinear form on $V_0$ is definite, the element $g$ in the assertion can be taken from $\LieGL(V_1)_o\times \LieSO(V_0)$.
	Note that, for any $S \in \lie{c}$, we have $S \in \lie{h}^\perp$ if and only if $S \in \RR(T_1 - T_2)$.
\end{proof}

\begin{lemma} \label{lem:SOpq2-degenerate}
	Let $T \in \lie{g}_1$ such that $\dim(\Im(T)) = 1$.
	Then there exists $g \in G$ such that $\Ad(g)T \in \RR T_3 \subset \lie{h}^\perp$.
\end{lemma}

\begin{proof}
	Since the bilinear form on $V_0$ is definite, the element $g$ in the assertion can be taken from $\LieGL(V_1)_o\times \LieSO(V_0)$.
\end{proof}

\begin{theorem} \label{thm:SOpq2}
	If $\lie{g} \simeq \lieSO(2, q)$ ($q \geq 5$) and $\lie{h} \simeq \lieSU(1, k)$ ($2\leq k\leq q/2$), then $(\lie{g}, \lie{h}, \orbit{O})$ is not minimal for any $\orbit{O}$.
\end{theorem}

\begin{proof}
	If $2k = q$, then $(\lie{g}, \lieNorm_{\lie{g}}(\lie{h}))$ is a symmetric pair and hence $(\lie{g}, \lie{h}, \orbit{O})$ is not minimal for any $\orbit{O}$ by Lemma \ref{lem:SufficientSymmetric}.
	If $\orbit{O}$ contains an element $T \in \lie{g}_1$ with $\rank(T) = 2$, the assertion follows from Lemmas \ref{lem:ReductionRealStronglyOrthogonal} and \ref{lem:SOpq2-nondegenerate}.
	If not, the assertion follows from Lemmas \ref{lem:SufficientMaxAbelian} and \ref{lem:SOpq2-degenerate}.
\end{proof}


\subsection{Minimality} \label{subsect:AssociatedPair}

Let $(\lie{g}, \lie{h})$ be a pair of type $BC_1$ with $\lie{h} \simeq \lieSU(1, k)$ ($k \geq 2$).
We shall give a necessary condition for minimality.

Fix $h \in \lie{h}^{-\theta}$ as in Definition \ref{def:typeBC}.
Let $\sigma$ be the involution of $\lie{g}$ such that $\lie{g}^\sigma = \lie{g}_{-2}\oplus\lie{g}_0 \oplus \lie{g}_2 \subset \lie{g}$.
Then $\lie{h}$ is $\sigma$-stable and we have $\lie{g}^{\theta\sigma} \supset \lie{g}^{-\theta, -\sigma} = (\lie{g}_{-1}\oplus \lie{g}_1) \cap \lie{g}^{-\theta}$.

Fix an element $Z \in \lie{g}_2$ such that $\set{h, Z, -\theta(Z)}$ forms an $\lie{sl}_2$-triple.
Then $Z + \theta(Z) \in \lie{k}^\sigma$.
Set $J\coloneq Z + \theta(Z)$.
Since $\dim_\RR(\lie{g}_2) = 1$, we have
\begin{align}
	\lie{k}^\sigma = ((\lie{g}_2 \oplus \lie{g}_{-2}) \cap \lie{k}) \oplus (\lie{g}_0\cap \lie{k}) = \RR J \oplus (\lie{k}\cap \lie{g}_0). \label{eqn:KsigmaDecomp}
\end{align}
This implies that $J \in \lieCent(\lie{k}^\sigma)$.

\begin{lemma}\label{lem:ComplexStructure}
	$\ad(J)$ gives a $\lie{k}^{\sigma}$-invariant complex structure on $\lie{g}_1\oplus \lie{g}_{-1}$.
\end{lemma}

\begin{proof}
	The invariance of $\ad(J)$ follows from $J \in \lieCent(\lie{k}^\sigma)$.
	Since $[Z, \lie{g}_1] = 0$ and $[\theta(Z), \lie{g}_{-1}] = 0$, we have
	\begin{align*}
		\ad(J)^2(X) &= [Z, [\theta(Z), X]] = [[Z, \theta(Z)], X] = [-h, X] = -X
	\end{align*}
	for any $X \in \lie{g}_1$.
	Similarly, we have $\ad(J)^2(Y) = -Y$ for any $Y \in \lie{g}_{-1}$.
	Hence $\ad(J)$ defines a complex structure on $\lie{g}_1 \oplus \lie{g}_{-1}$.
\end{proof}

Take a maximal torus $\lie{t}'$ in $(\lie{k}_H\cap \lie{g}_0) \oplus \lieCent_{\lie{g}}(\lie{h}) \subset \lie{k}^\sigma$.
Then $\lie{t}\coloneq \lie{t}'\oplus \RR J$ is a compact Cartan subalgebra of $\lie{g}$ by Lemma \ref{lem:EmbeddingSU}, and $\lie{t}$ is contained in $\lie{k}^{\sigma} \subset \lie{g}^{\theta\sigma}$.
Set $\lie{p} \coloneq \lie{g}^{-\theta, -\sigma}_{\CC}$ and let $\lie{p} = \lie{p}_+\oplus \lie{p}_-$ be the eigenspace decomposition of $\lie{p}$ with respect to the action of $\sqrt{-1}(Z+\theta(Z))$, where $\lie{p}_+$ and $\lie{p}_-$ are the eigenspaces with eigenvalues $1$ and $-1$, respectively.
Take a set of strongly orthogonal roots $\set{\gamma_1, \ldots, \gamma_r}$ in $\Delta(\lie{p}_+\cap (\lieC{h})^\perp, \lieC{t})$ with $r$ as large as possible.
Take root vectors $X_{\gamma_i} \in (\lieC{g})_{\gamma_i}$ for $1 \leq i \leq r$.

Set $\lie{c} \coloneq \RR h \oplus \lie{t'}$ and, for $1\leq i\leq r$, define $\gamma^{\pm}_i \in \lie{c}^*$ by
\begin{align*}
	\gamma^{+}_i(h) &= 1, \quad \gamma^{+}_i|_{\lie{t'}} = \gamma_i|_{\lie{t'}}, \\
	\gamma^{-}_i(h) &= -1, \quad \gamma^{-}_i|_{\lie{t'}} = \gamma_i|_{\lie{t'}}.
\end{align*}
Then we have
\begin{align*}
	\overline{(\lieC{g})_{\gamma^{\pm}_i}} = (\lieC{g})_{-\gamma^{\mp}_i}, \quad \theta((\lieC{g})_{\gamma^{\pm}_i}) = (\lieC{g})_{\gamma^{\mp}_i}.
\end{align*}
Let $\lieC{g}'$ be the subalgebra of $\lieC{g}$ generated by
\begin{align*}
	(\lieC{g})_{\gamma^{\pm}_i}, (\lieC{g})_{-\gamma^{\pm}_i} \ (1\leq i \leq r) \text{ and } \lieC{c},
\end{align*}
and set $\lie{g'} \coloneq \lie{g}\cap \lieC{g}'$, which is a real form of $\lieC{g}'$.
Then $\lie{g'}$ is a $\theta$-stable subalgebra of $\lie{g}$ containing $h$.

We shall discuss when $\lie{g'}$ satisfies the assumptions of Proposition \ref{prop:NonMinimalSubalgebra}.
Set
\begin{align*}
	\lie{a}' \coloneq \bigoplus_{i=1}^r \RR(X_{\gamma_i} + \overline{X_{\gamma_i}}) \subset \lie{g}^{-\theta, -\sigma}.
\end{align*}
Then $\lie{a'}$ is an abelian subspace of $\lie{g}^{-\theta, -\sigma}$, and we have
\begin{align}
	\lie{a'} \subset \lie{g'} \cap \lie{h}^\perp, \quad \ad(J)\lie{a'} = \bigoplus_{i=1}^r \RR\sqrt{-1}(X_{\gamma_i} - \overline{X_{\gamma_i}}) \subset \lie{g'} \cap \lie{h}^\perp. \label{eqn:SmallSubalgebraA}
\end{align}

\begin{lemma} \label{lem:SmallSubalgebraStronglyOrthogonal}
	$\lie{g'}$ satisfies the following properties.
	\begin{enumparen}
		\item $X_{\gamma_i} + \overline{X_{\gamma_i}}, \sqrt{-1}(X_{\gamma_i} - \overline{X_{\gamma_i}}) \in \lie{g'}$.
		\item Let $\lie{g}'_{ss}$ be the semisimple part of $\lie{g'}$. Then one has
		\begin{align*}
			r \leq \rank_{\RR}(\lie{g}'_{ss}) \leq \rank(\lie{g}'_{ss}) \leq r + 1.
		\end{align*}
		\item $\lie{g'} = (\lie{g'} \cap \lie{h}) \oplus (\lie{g'} \cap \lie{h}^\perp)$.
	\end{enumparen}
\end{lemma}

\begin{proof}
	The assertion (1) is clear from the definition of $\lie{g'}$.
	By definition, the weight $\gamma^+_i - \gamma^-_i$ does not depend on $i$.
	Hence we have
	\begin{align*}
		r = \dim_{\RR}(\lie{a'}) &\leq \rank_{\RR}(\lie{g}'_{ss}) \leq \rank(\lie{g}'_{ss}) \\
		&= \dim_{\RR}\mathrm{span}_{\RR}\set{\gamma^+_1, \ldots, \gamma^+_r, \gamma^-_1, \ldots, \gamma^-_r} \leq r + 1.
	\end{align*}
	This shows (2).

	Since $\lie{c}$ is a Cartan subalgebra contained in $\lieNorm_{\lie{g}}(\lie{h})$, any root space with respect to $\lieC{c}$ is contained in $\lieC{h}$ or $\lieC{h}^\perp$.
	This proves assertion (3).
\end{proof}

\begin{corollary} \label{cor:SmallG'}
	If $\lie{g'} = \lie{g}$, then one has
	\begin{align*}
		r \leq \rank_{\RR}(\lie{g}) \leq \rank(\lie{g}) \leq r + 1.
	\end{align*}
\end{corollary}

\begin{lemma} \label{lem:SufficientForSecondReduction}
	Let $\orbit{O}$ be a nilpotent orbit in $\lie{g}$ with $\overline{\orbit{O}} \cap \Nilpotent(\lie{h}^\perp, H) \neq \set{0}$.
	Assume that there exists a compact subgroup $K_0$ of $C_K(h)$ such that $\lie{g}_1$ is $K_0$-stable and $\exp(\RR \ad(J))\Ad(K_0) \lie{a'} \supset \lie{g}^{-\theta, -\sigma} \cap \lie{h}^\perp$.
	Then one has $\overline{\orbit{O}} \cap \lie{h}^\perp \cap \lie{g}'_1 \neq \set{0}$.
\end{lemma}

\begin{proof}
	By Lemma \ref{lem:ExistenceWeightVector}, there exists a non-zero element $X \in \overline{\orbit{O}} \cap \lie{g}_1$.
	Then, we have $X - \theta(X) \in \lie{g}^{-\theta, -\sigma} \cap \lie{h}^\perp$.
	By assumption, there exist $k \in K_0$ and $s \in \RR$ such that $\exp(s\ad(J))\Ad(k)(X - \theta(X)) \in \lie{a'}$.
	By \eqref{eqn:SmallSubalgebraA}, we have
	\begin{align*}
		\Ad(k)X &\in \exp(\RR \ad(J))\lie{a'} + \ad(h)\exp(\RR \ad(J))\lie{a'} \\
		&\subset \lie{a'} + \ad(J)(\lie{a'}) + \ad(h)\lie{a'} + \ad(h)\ad(J)(\lie{a'}) \subset \lie{g'} \cap \lie{h}^\perp.
	\end{align*}
	This implies that $\overline{\orbit{O}} \cap \lie{h}^\perp \cap \lie{g}'_1 \neq \set{0}$.
\end{proof}

\begin{proposition} \label{prop:NonMinimalSymmetric}
	Suppose that $\lie{g} \neq \lie{g'}$ and $(\lie{g}^{\theta\sigma}, \lieNorm_{\lie{g}}(\lie{h})^{\theta\sigma})$ is a symmetric pair.
	Then $(\lie{g}, \lie{h}, \orbit{O})$ is not minimal for any $\orbit{O}$.
\end{proposition}

\begin{proof}
	Let $\tau$ be an involution of $\lie{g}^{\theta\sigma}$ corresponding to $\lieNorm_{\lie{g}}(\lie{h})^{\theta\sigma}$.
	Note that $\lieNorm_{\lie{g}}(\lie{h})^{\theta\sigma}$ contains $\lie{t}$ and the center of $\lie{g}^{\theta\sigma}$.
	Then $\tau$ is unique, $\tau$ commutes with $\theta$ and $\sigma$, and the restriction of the form $(\cdot, \cdot)$ to $\lie{g}^{\theta\sigma}$ is $\tau$-invariant.

	Set $\lie{s} \coloneq \lie{g}^{\theta\sigma, \theta\tau}$.
	Since $\lieCent_{\lie{g}}(\lie{h})$ is compact, we have
	\begin{align*}
		\lie{s}^{-\theta} &= \lie{g}^{-\theta, -\sigma, -\tau} = \lie{g}^{\theta\sigma} \cap \lie{g}^{-\theta} \cap \lie{h}^\perp, \\
		\lie{t}&\subset \lie{g}^{\theta, \sigma, \tau} \subset \lie{s}.
	\end{align*}
	Hence $\lie{s}^{-\theta}$ has an $\lie{s}^{\theta}$-invariant complex structure and we have $\lie{a}' \subset \lie{s}^{-\theta}$.
	In other words, any non-compact simple factor of $\lie{s}$ is of Hermitian type.
	Since $\lie{a'}$ is constructed from a maximal set of strongly orthogonal roots, $\lie{a'}$ is a maximal abelian subspace of $\lie{s}^{-\theta}$ by \cite[Lemma 7.143]{Kn02}.

	Let $K_0$ be the analytic subgroup of $G$ with the Lie algebra $\lieNorm_{\lie{k}}(\lie{h})\cap \lie{g}_0 \subset \lie{g}^{\theta, \sigma, \tau}$, and set $K'\coloneq \exp(\RR J)K_0$.
	By \eqref{eqn:KsigmaDecomp}, we have
	\begin{align*}
		\lie{g}^{\theta, \sigma, \tau} = (\RR (Z+\theta(Z)) \oplus (\lie{k}\cap \lie{g}_0)) \cap \lieNorm_{\lie{g}}(\lie{h})
		= \RR J \oplus (\lieNorm_{\lie{k}}(\lie{h})\cap \lie{g}_0).
	\end{align*}
	Then $\lie{g}_1$ is $K_0$-stable and $K'$ has the Lie algebra $\lie{g}^{\theta, \sigma, \tau}$.
	Since $\lie{a'}$ is a maximal abelian subspace of $\lie{s}^{-\theta}$, we have
	\begin{align*}
		\Ad(K')\lie{a'} = \lie{s}^{-\theta} = \lie{g}^{-\theta, -\sigma} \cap \lie{h}^\perp.
	\end{align*}
	Therefore the assertion follows from Lemma \ref{lem:SufficientForSecondReduction}.
\end{proof}

\begin{proposition} \label{prop:NonMinimalFullRank}
	Suppose that $\lie{g} \neq \lie{g'}$ and $r = \rank_{\RR}(\lie{g}^{\theta\sigma})$.
	Then $(\lie{g}, \lie{h}, \orbit{O})$ is not minimal for any $\orbit{O}$.
\end{proposition}

\begin{proof}
	Let $K_0$ be the analytic subgroup of $G$ with the Lie algebra $\lie{k}\cap \lie{g}_0$, and set $K'\coloneq \exp(\RR J)K_0$.
	Then $\lie{g}_1$ is $K_0$-stable and $K'$ has the Lie algebra $\lie{k}^\sigma$ by \eqref{eqn:KsigmaDecomp}.
	By assumption, $\lie{a'}$ is a maximal abelian subspace of $\lie{g}^{-\theta, -\sigma}$.
	Hence we have
	\begin{align*}
		\Ad(K')\lie{a'} = \lie{g}^{-\theta, -\sigma} \supset \lie{g}^{-\theta, -\sigma} \cap \lie{h}^\perp.
	\end{align*}
	Therefore the assertion follows from Lemma \ref{lem:SufficientForSecondReduction}.
\end{proof}

The conditions in Propositions \ref{prop:NonMinimalSymmetric} and \ref{prop:NonMinimalFullRank} can be verified purely in terms of root systems.
These propositions are closely related to Proposition \ref{prop:ReductionRealStronglyOrthogonal}, but they require no ideal decomposition on $\lie{g}^{\theta\sigma}$ and therefore apply in a different range of cases.

\begin{lemma} \label{lem:NonMinimalSOstar}
	Suppose that $\lie{g} \simeq \lieSO^*(2n)$ ($n \geq 5$).
	Then there exists no minimal NDD tuple $(\lie{g}, \lie{h}, \orbit{O})$.
\end{lemma}

\begin{proof}
	If $(\lie{g}, \lieNorm_{\lie{g}}(\lie{h}))$ is a symmetric pair, then $(\lie{g}, \lie{h}, \orbit{O})$ is not minimal for any $\orbit{O}$ by Lemma \ref{lem:SufficientSymmetric}.
	By Lemma \ref{lem:EmbeddingSUtoSOstar}, we may assume that $\lieNorm_{\lie{g}}(\lie{h}) \simeq \lieU(1, n-2) \oplus \lieSO^*(2)$.
	We shall describe the embedding explicitly.

	Let $F = F_+ \oplus F_- \oplus F_0$ be a complex vector space with a Hermitian form $B$ of signature $(2, n - 2)$ whose restrictions to each of $F_+$ and $F_-$ have signature $(1, 0)$ and whose restriction to $F_0$ has signature $(0, n - 2)$.
	Identify $\lie{g}$ with the Lie algebra of the isometry group of $\HH\otimes_{\CC} F$.
	Fix a subspace $W_0\subset F_0$ such that $\dim(W_0) = n - 3$.
	Set $W'_0 \coloneq W_0^\perp \cap F_0$.

	We may choose the embedding $\lie{h} \hookrightarrow \lie{g}$ and $h \in \lie{h}$ such that
	\begin{align*}
		\lie{g}^{\sigma} &= \lieSO^*(\HH\otimes_{\CC} (F_{+}\oplus F_{-})) \oplus \lieSO^*(\HH\otimes_{\CC}F_0), \\
		\lie{g}^{\theta\sigma} &= \lieU(F), \quad \lie{g}^{\theta, \sigma} = \lieU(F_+ \oplus F_-) \oplus \lieU(F_0), \\
		\lie{h} &= \lieSU(F_{+} \oplus jF_- \oplus W_0).
	\end{align*}
	Then we have
	\begin{align*}
		\lieNorm_{\lie{g}}(\lie{h}) &= \lieU(F_{+} \oplus jF_-\oplus W_0) \oplus \lieSO^*(\HH\otimes_{\CC} W'_0), \\
		\lieNorm_{\lie{g}}(\lie{h})^{\theta\sigma} &= \lieU(F_{+} \oplus W_0) \oplus \lieU(F_-) \oplus \lieU(W'_0).
	\end{align*}
	$\lieNorm_{\lie{g}}(\lie{h})^{\theta\sigma}$ is a Levi subalgebra of $\lie{g}^{\theta\sigma}$.
	Hence we may take a coordinate system $\set{\varepsilon_1, \ldots, \varepsilon_n}$ of $\lieC{t}$ such that
	\begin{align*}
		\Delta(\lieC{g}^{\theta\sigma}, \lieC{t}) &= \set{\pm (\varepsilon_i - \varepsilon_j) : 1\leq i < j \leq n}, \\
		\Delta(\lie{p}_+\cap \lieC{h}, \lieC{t}) &= \set{\varepsilon_1 - \varepsilon_i : 3\leq i \leq n-1}, \\
		\Delta(\lie{p}_+\cap (\lieC{h})^\perp, \lieC{t}) &= \set{ \varepsilon_2 - \varepsilon_i : 3 \leq i \leq n-1} \cup \set{\varepsilon_1 - \varepsilon_n, \varepsilon_2 - \varepsilon_n}.
	\end{align*}
	Since $\set{\varepsilon_2 - \varepsilon_3, \varepsilon_1 - \varepsilon_n}$ is a maximal subset of strongly orthogonal roots in $\Delta(\lie{p}_+\cap (\lieC{h})^\perp, \lieC{t})$, we have $r = 2 = \rank_{\RR}(\lie{g}^{\theta\sigma}) < \rank(\lie{g}) - 1$.
	By Proposition \ref{prop:NonMinimalFullRank}, the assertion follows.
\end{proof}

Summarizing Lemmas \ref{lem:EmbeddingSUtoSp}, \ref{lem:EmbeddingSUtoSU}, \ref{lem:NonMinimalSOstar} and Theorems \ref{thm:SOpq4}, \ref{thm:SOpq3} and \ref{thm:SOpq2}, we obtain the following theorem.

\begin{theorem} \label{thm:NonMinimalClassical}
	If $\lie{g}$ is a classical simple Lie algebra and $\lie{h} \simeq \lieSU(1, k)$ ($k \geq 2$), then there exists no minimal NDD tuple $(\lie{g}, \lie{h}, \orbit{O})$.
\end{theorem}

\subsection{Embeddings of \texorpdfstring{$\lieSU(1, k)$}{su(1,k)}: Exceptional \texorpdfstring{$\lie{g}$}{g}}

We shall consider the case where $\lie{g}$ is an exceptional Lie algebra.
We classify the embeddings of $\lieSU(1, k)$ into $\lie{g}$ for which the pair $(\lie{g}, \lie{h})$ is of type $BC_1$ (Definition \ref{def:typeBC}).
The classification is obtained by implementing the following straightforward algorithm on a computer.
The source code and the resulting classification data are available in \cite{Ki26-code}.
Our classification is up to isomorphism of pairs $(\lie{g}, \lie{h})$.

Before giving the algorithm, we shall state some observations on the embedding of $\lie{h} \simeq \lieSU(1, k)$ into $\lie{g}$.
Let $(\lie{g}, \lie{h})$ be a pair of type $BC_1$ and $\lie{h} \simeq \lieSU(1, k)$.
Let $h \in \lie{h}^{-\theta}$ be as in Definition \ref{def:typeBC} and $\sigma$ be the involution of $\lie{g}$ such that $\lie{g}^\sigma = \lie{g}_{-2} \oplus \lie{g}_0 \oplus \lie{g}_2$.
Take a $\sigma$-stable maximal torus $\lie{t}$ of $\lieNorm_{\lie{g}}(\lie{h})^\theta$.
Then we have
\begin{align*}
	\lie{t} = (\lie{t}\cap (\lie{g}_{-2} \oplus \lie{g}_2)) \oplus (\lie{t}\cap \lie{g}_0),
\end{align*}
and $\lie{t}$ is a compact Cartan subalgebra of $\lie{g}$ by Lemma \ref{lem:EmbeddingSU}.
Note that
\begin{align*}
	\lieNorm_{\lie{g}}(\lie{h})^\sigma = (\lie{g}_{-2}\oplus \RR h \oplus \lie{g}_2) \oplus (\lie{k}_H\cap \lie{g}_0) \oplus \lieCent_{\lie{g}}(\lie{h}),
\end{align*}
where the first summand is isomorphic to $\lieSU(1, 1) \simeq \lie{sl}(2, \RR)$.

Each root in $\Delta(\lieC{g}, \lieC{t})$ is either compact or non-compact.
Note that one of the Vogan diagrams of $\lie{h} \simeq \lieSU(1, k)$ is given by the following diagram.
\begin{align}
	\dynkin[scale=2, labels={\alpha_1, \alpha_2, \alpha_3, \alpha_{k-1}, \alpha_k}]{A}{*oo.oo} \label{eqn:VoganSU}
\end{align}
In the following, we use the same labeling of simple roots for $\lie{h}$ as in the above diagram.

\begin{lemma} \label{lem:SimpleRootEmbeddingSU}
	Let $S \coloneq \set{\alpha_1, \ldots, \alpha_k}$ be a set of simple roots in $\Delta(\lieC{h}, \lieC{t})$ with the Vogan diagram \eqref{eqn:VoganSU}.
	\begin{enumparen}
		\item $S$ is a set of simple roots of a root subsystem of $\Delta(\lieC{g}, \lieC{t})$ of type $A_k$.
		\item $\alpha_1$ is the only non-compact root in $S$.
		\item Any root in $\Delta(\lieC{g}, \lieC{t})$ that is strongly orthogonal to all roots in $S$ is compact.
		\item Any root in $S$ (and hence any root in $\Delta(\lieC{h}, \lieC{t})$) is long in $\Delta(\lieC{g}, \lieC{t})$.
	\end{enumparen}
\end{lemma}

\begin{proof}
	(1) and (2) are clear.
	Since $\lieCent_{\lie{g}}(\lie{h})$ is a compact subalgebra of $\lie{g}$, (3) is satisfied.

	Set $\lie{t'} \coloneq \lie{t} \cap \lie{g}_0$.
	Then $\lie{t'}$ is a maximal torus in $(\lie{g}_0 \cap \lie{k}_H) \oplus \lieCent_{\lie{g}}(\lie{h})$, and $\lie{c} \coloneq \RR h \oplus \lie{t'}$ is a Cartan subalgebra of $\lie{g}$.
	By Lemma \ref{lem:G2RootSpace}, $(\lieC{g})_2$ is a root space of $\lieC{g}$ with respect to $\lieC{c}$.
	Let $\alpha' \in \lieC{c}^*$ be the root such that $(\lieC{g})_2 = (\lieC{g})_{\alpha'}$.
	Then $\alpha'$ is a real root and $2(\alpha', \beta)/(\alpha', \alpha') = \beta(h) \in \set{\pm 1, 0}$ for any $\beta \in \Delta(\lieC{g}, \lieC{c})\backslash\set{\pm \alpha'}$.
	Hence $\alpha'$ is a long root.

	Let $\psi$ be the Cayley transform (see \cite[Chapter VI.7]{Kn02}) by the real root $\alpha'$.
	Then we have $\lie{t} =\psi(\lieC{c}) \cap \lie{g}$ and $\alpha'_1\coloneq \alpha' \circ \psi^{-1}$ is a non-compact long root in $\Delta(\lieC{g}, \lieC{t})$.
	Since $\Delta(\lieC{h}, \lieC{t})$ is of type $A_k$, any root has the same length.
	Hence any root in $\Delta(\lieC{h}, \lieC{t})$ is long in $\Delta(\lieC{g}, \lieC{t})$.
	This shows (4).
\end{proof}

We shall consider the converse of Lemma \ref{lem:SimpleRootEmbeddingSU}.
For a while, let $\lie{g}$ be a simple real Lie algebra with a compact Cartan subalgebra $\lie{t} \subset \lie{g}^\theta$.

\begin{lemma} \label{lem:SimpleRootEmbeddingSUInv}
	Let $S = \set{\alpha_1, \ldots, \alpha_k}$ be a subset of $\Delta(\lieC{g}, \lieC{t})$ satisfying the conclusions of Lemma \ref{lem:SimpleRootEmbeddingSU}.
	Let $\lie{h}$ be the subalgebra whose root system $\Delta(\lieC{h}, \lieC{t})$ has $S$ as a set of simple roots.
	Then $(\lie{g}, \lie{h})$ is a pair of type $BC_1$.
\end{lemma}

\begin{proof}
	We shall verify the conditions (1)--(5) in Definition \ref{def:typeBC}.
	The conditions (1) and (2) are clear.
	By assumption (3) (in Lemma \ref{lem:SimpleRootEmbeddingSU}), $\lieCent_{\lie{g}}(\lie{h})$ is a compact subalgebra of $\lie{g}$.
	This is (3) in Definition \ref{def:typeBC}.

	Let $\varphi \in \Aut(\lieC{g})$ be the Cayley transform (see \cite[Chapter VI.7]{Kn02}) by the non-compact root $\alpha_1$, and set $\lie{c} \coloneq \varphi(\lieC{t}) \cap \lie{g}$.
	Then $\varphi$ normalizes $\lieC{h}$ and $\alpha'_1\coloneq \alpha_1 \circ \varphi^{-1}$ is a unique real root in $\Delta(\lieC{g}, \lieC{c})$ up to sign.
	Since $\alpha'_1$ is a long root, $2(\alpha'_1, \beta')/(\alpha'_1, \alpha'_1) \in \set{\pm 2, \pm 1, 0}$ for any $\beta' \in \Delta(\lieC{g}, \lieC{c})$, and $2(\alpha'_1, \beta')/(\alpha'_1, \alpha'_1) = 2$ only if $\beta' = \alpha'_1$.
	This shows the conditions (4) and (5).
\end{proof}

Our task is to classify the subsets $S\coloneq \set{\alpha_1, \ldots, \alpha_k} \subset \Delta(\lieC{g}, \lieC{t})$ satisfying the conditions in Lemma \ref{lem:SimpleRootEmbeddingSU}.

\begin{algorithm} \label{alg:EmbeddingSU}
	Find and fix a non-compact root $\alpha_1$ in $\Delta(\lieC{g}, \lieC{t})$.
	If there are two lengths of roots, choose $\alpha_1$ to be a long root.
	Let $\Delta = \set{\beta_1, \ldots, \beta_m}$ be an ordered set of compact (long) roots in $\Delta(\lieC{g}, \lieC{t})$.
	Consider $\Delta \cup \set{\alpha_1}$ as an undirected graph in which two vertices $\alpha$ and $\beta$ are joined by an edge labeled $(\alpha, \beta)$ whenever $(\alpha, \beta) \neq 0$.
	Then $S$ is identified with a path that starts at $\alpha_1$ and has negatively labeled edges; moreover, the vertices $\alpha_i$ and $\alpha_j$ are nonadjacent whenever $|i-j|>1$.
	Enumerate all such paths $S = \set{\alpha_1, \ldots, \alpha_k}$ ($k\geq 2$) by depth-first search, and verify that no non-compact root is strongly orthogonal to every root in $S$.
\end{algorithm}

\begin{theorem}
	Algorithm \ref{alg:EmbeddingSU} produces a list of embeddings of $\lieSU(1, k)$ into the simple real Lie algebra $\lie{g}$ for which the corresponding pairs satisfy the conditions in Definition \ref{def:typeBC}.
	The list may contain mutually isomorphic entries, but every such embedding is isomorphic to one appearing in the list.
\end{theorem}

\begin{proof}
	Lemma \ref{lem:SimpleRootEmbeddingSUInv} essentially completes the proof.
	It remains only to justify the constraint that $\alpha_1$ is fixed in Algorithm \ref{alg:EmbeddingSU}, and this follows from the following fact.
	By \cite[Chapter XII. Problem 21]{Kn01}, any two non-compact long roots are conjugate by the group
	\begin{align*}
		\set{g \in \Aut(\Delta(\lieC{g}, \lieC{t})): g\text{ preserves compact roots}}.
	\end{align*}
	This result follows from the fact that $\lieC{g}^{-\theta}$ is irreducible or the direct sum of two irreducible representations of $\lieC{k}$.
	This and Lemmas \ref{lem:SimpleRootEmbeddingSU} and \ref{lem:SimpleRootEmbeddingSUInv} prove the theorem.
\end{proof}

Algorithm \ref{alg:EmbeddingSU} often produces mutually isomorphic entries.
Therefore, from the resulting list, we need to choose a set of representatives under the action of the automorphism group $\Aut(\Delta(\lieC{g}, \lieC{t}))$ preserving compactness.
For simplicity, in our implementation, we first choose representatives under the Weyl group of the root system $\set{\alpha \in \Delta(\lieC{k}, \lieC{t}): (\alpha, \alpha_1) = 0}$ and then eliminate the remaining duplicates by comparing their Vogan diagrams.

In the following, we give the list of embeddings of $\lieSU(1, k)$.
Retain the notation $\lie{g}, \lie{h}, \sigma, \lie{t}$ from the beginning of this subsection.
As in Subsections \ref{subsect:ReductionStronglyOrthogonal} and \ref{subsect:AssociatedPair}, the pair $(\lie{g}^{\theta\sigma}, \lieNorm_{\lie{g}}(\lie{h})^{\theta\sigma})$ plays an important role in establishing the non-minimality of NDD tuples.
We describe the embeddings using Vogan diagrams of $\lie{g}$ and $\lie{g}^{\theta\sigma}$.
Fix a set of simple roots $\set{\beta_1, \ldots, \beta_l}$ of $\Delta(\lieCent_{\lieC{g}}(\lieC{h}), \lieC{t})$.
Then $\set{\alpha_1, \ldots, \alpha_k, \beta_1, \ldots, \beta_l}$ is a set of simple roots of $\Delta(\lieNorm_{\lieC{g}}(\lieC{h}), \lieC{t})$.
The classification of the pairs $(\lie{g},\lie{h})$ below yields the following propositions.

\begin{proposition}
	One has $k+l = \rank(\lie{g})$ or $\rank(\lie{g}) - 1$.
	If $k+l=\rank(\lie{g}) - 1$, then $\lieNorm_{\lieC{g}}(\lieC{h})$ is a Levi subalgebra of $\lieC{g}$ with one-dimensional center.
\end{proposition}
	
There are several ways to realize the Vogan diagram of $\lieNorm_{\lie{g}}(\lie{h})$ in that of $\lie{g}$.
We use the following convention.

\begin{proposition}
	Suppose that $\lie{g}$ is exceptional.
	There exists a unique positive system of $\Delta(\lieC{g}, \lieC{t})$ such that $\alpha_1$ is a lowest root and $\alpha_2, \ldots, \alpha_k, \beta_1, \ldots, \beta_l$ are simple roots.
\end{proposition}

\begin{example}
	We give a concrete example in which $\lie{g}$ is classical rather than exceptional.
	Let $(\lie{g}, \lie{h}) = (\lieSU(4,3), \lieSU(1,3))$.
	Then we have $\lieNorm_{\lie{g}}(\lie{h}) = \lie{s}(\lieU(1,3) \oplus \lieU(3, 0))$.
	The embedding is represented by the following Vogan diagram with an extended vertex $\alpha_1$.
	\begin{align*}
		\dynkin[scale=2, extended, labels={\alpha_1, \alpha_2, \alpha_3, , \beta_1, \beta_2, }, affine mark=*]{A}{oo*ooo}
	\end{align*}
\end{example}

To realize the Vogan diagram of $\lieNorm_{\lie{g}}(\lie{h})^{\theta\sigma}$ in that of $\lie{g}^{\theta\sigma}$, we consider the set of simple roots $\set{\alpha_1 + \alpha_2, \alpha_3, \ldots, \alpha_k}$ of $\Delta((\lieC{h})^{\theta\sigma}, \lieC{t})$.
Here $\alpha_1 + \alpha_2$ is a non-compact root.
Note that $\lie{h}^{\theta\sigma}$ is isomorphic to $\lie{s}(\lieU(1, k-1)\oplus \lieU(1))$.

\begin{proposition} \label{prop:EmbeddingVoganAssociatedPair}
	Suppose that $\lie{g}$ is exceptional.
	There exists a positive system of $\Delta((\lieC{g})^{\theta\sigma}, \lieC{t})$ such that $\alpha_1 + \alpha_2, \alpha_3, \ldots, \alpha_k, \beta_1, \ldots, \beta_l$ are simple roots.
\end{proposition}

\begin{example}
	Let $(\lie{g}, \lie{h}) = (\lieSU(4,3), \lieSU(1,3))$.
	Then we have $\lie{g}^{\theta\sigma} = \lie{s}(\lieU(1,2) \oplus \lieU(3, 1))$ and $\lieNorm_{\lie{g}}(\lie{h})^{\theta\sigma} = \lie{s}(\lieU(1, 2) \oplus \lieU(0, 1) \oplus \lieU(3, 0))$.
	The embedding is represented by the following Vogan diagram.
	\begin{align*}
		\dynkin[scale=2, labels={\alpha_1 + \alpha_2, \alpha_3}]{A}{*o}\quad 
		\dynkin[scale=2, labels={, \beta_1, \beta_2}]{A}{*oo}
	\end{align*}
\end{example}

In Proposition \ref{prop:EmbeddingVoganAssociatedPair}, the positive system may not be unique.
We choose one of them to minimize the number of non-compact roots and realize the Vogan diagram of $\lieNorm_{\lie{g}}(\lie{h})^{\theta\sigma}$ in that of $\lie{g}^{\theta\sigma}$.

\begin{proposition} \label{prop:EmbeddingVoganAssociatedPair2}
	Among the Vogan diagrams provided by Proposition \ref{prop:EmbeddingVoganAssociatedPair}, there exists a unique Vogan diagram of $\Delta((\lieC{g})^{\theta\sigma}, \lieC{t})$ in which each connected component has a unique non-compact simple root.
\end{proposition}

Note that the Vogan diagrams above have information on the semisimple parts of $\lie{g}^{\theta\sigma}$, $\lieNorm_{\lie{g}}(\lie{h})$ and $\lieNorm_{\lie{g}}(\lie{h})^{\theta\sigma}$.
However, since they contain the compact Cartan subalgebra $\lie{t}$, we can recover the three subalgebras from the Vogan diagrams.

\begin{theorem} \label{thm:NonMinimalExceptional}
	If $(\lie{g}, \lie{h})$ is a pair of type $BC_1$ with $\lie{h} \simeq \lieSU(1, k)$ and with exceptional $\lie{g}$, then $(\lie{g}, \lie{h}, \orbit{O})$ is not minimal for any $\orbit{O}$.
\end{theorem}

We shall show Theorem \ref{thm:NonMinimalExceptional} using Propositions \ref{prop:ReductionRealStronglyOrthogonal}, \ref{prop:NonMinimalSymmetric} and \ref{prop:NonMinimalFullRank}.
We use the same notation $r$ and $\lie{p}_+$ as in Subsection \ref{subsect:AssociatedPair}.

\subsubsection{\texorpdfstring{$\lie{g} = \lieE_{6(2)}$}{g=e6(2)}}

% quaternionic

We have $\lie{g}^{\theta\sigma} \simeq \lieSU(3, 3) \oplus \lieSL(2, \RR)$.
Then $r$ in Subsection \ref{subsect:AssociatedPair} satisfies
\begin{align*}
	r \leq \rank_{\RR}(\lie{g}^{\theta\sigma}) = 4 < 5 = \rank(\lie{g}) - 1.
\end{align*}
There are two embeddings of $\lieSU(1, k)$ into $\lieE_{6(2)}$.

\begin{align*}
	(\lie{g}, \lieNorm_{\lie{g}}(\lie{h})): \quad & \dynkin[scale=2, extended, labels={\alpha_1, \beta_1, \alpha_2, \beta_2, , \beta_3, \beta_4}, affine mark=*]{E}{ooo*oo} \\
	(\lie{g}^{\theta\sigma}, \lieNorm_{\lie{g}}(\lie{h})^{\theta\sigma}): \quad & \dynkin[scale=2, labels={\beta_1, \beta_2, , \beta_3, \beta_4}]{A}{oo*oo} \dynkin[scale=2, labels={\alpha_2 + \alpha_1}]{A}{*}
\end{align*}

In the first embedding, we have $(\lie{g}^{\theta\sigma}, \lieNorm_{\lie{g}}(\lie{h})^{\theta\sigma}) \simeq (\lieSU(3, 3) \oplus \lieSL(2, \RR), \lie{s}(\lieU(3)\oplus \lieU(3)) \oplus \lieSL(2, \RR))$, which is a symmetric pair.
The assumptions of Proposition \ref{prop:NonMinimalSymmetric} hold.

\begin{align*}
	(\lie{g}, \lieNorm_{\lie{g}}(\lie{h})): \quad & \dynkin[scale=2, extended, labels={\alpha_1, , \alpha_2, \alpha_4, \alpha_3, , \beta_1}, affine mark=*]{E}{*ooo*o} \\
	(\lie{g}^{\theta\sigma}, \lieNorm_{\lie{g}}(\lie{h})^{\theta\sigma}): \quad & \dynkin[scale=2, labels={\alpha_4, \alpha_3, \alpha_2 + \alpha_1, , \beta_1}]{A}{oo*oo} \quad \dynkin[scale=2, labels={}]{A}{*}
\end{align*}

In the second embedding, we have $(\lie{g}^{\theta\sigma}, \lieNorm_{\lie{g}}(\lie{h})^{\theta\sigma}) \simeq (\lieSU(3, 3) \oplus \lieSL(2, \RR), \lie{s}(\lieU(3, 1)\oplus \lieU(2)) \oplus \lieSO(2))$, which is a symmetric pair.
The assumptions of Proposition \ref{prop:NonMinimalSymmetric} hold.

\subsubsection{\texorpdfstring{$\lie{g} = \lieE_{6(-14)}$}{g=e6(-14)}}

% Hermitian

We have $\lie{g}^{\theta\sigma} \simeq \lieSO^*(10) \oplus \lieSO(2)$.
Then $r$ in Subsection \ref{subsect:AssociatedPair} satisfies
\begin{align*}
	r \leq \rank_{\RR}(\lie{g}^{\theta\sigma}) = 2 < 5 = \rank(\lie{g}) - 1.
\end{align*}
There is one embedding of $\lieSU(1, k)$ into $\lieE_{6(-14)}$.

\begin{align*}
	(\lie{g}, \lieNorm_{\lie{g}}(\lie{h})): \quad & \dynkin[scale=2, extended, labels={\alpha_1, , \alpha_2, \alpha_4, \alpha_3, , \beta_1}, affine mark=*]{E}{*ooooo} \\
	(\lie{g}^{\theta\sigma}, \lieNorm_{\lie{g}}(\lie{h})^{\theta\sigma}): \quad & \dynkin[scale=2, labels={\beta_1, , \alpha_3, \alpha_4, \alpha_2 + \alpha_1}, label directions={,,right}]{D}{oooo*}
\end{align*}

In the embedding, we have $(\lie{g}^{\theta\sigma}, \lieNorm_{\lie{g}}(\lie{h})^{\theta\sigma}) \simeq (\lieSO^*(10) \oplus \lieSO(2), \lieSO^*(6)\oplus \lieU(2) \oplus \lieSO(2))$.
To take a set of strongly orthogonal roots, we fix a coordinate system $\set{\varepsilon_1, \ldots, \varepsilon_6}$ of $\lieC{t}$ such that
\begin{align*}
	\Delta(\lieC{g}^{\theta\sigma}, \lieC{t}) &= \set{\pm \varepsilon_i \pm \varepsilon_j: 1\leq i < j \leq 5}, \\
	\Delta(\lieC{g}^{\theta, \sigma}, \lieC{t}) &= \set{\pm (\varepsilon_i - \varepsilon_j): 1\leq i < j \leq 5}, \\
	\beta_1 = \varepsilon_1 - \varepsilon_2, &\quad \alpha_3 = \varepsilon_3 - \varepsilon_4, \quad \alpha_4 = \varepsilon_4 - \varepsilon_5, \quad \alpha_2 + \alpha_1 = \varepsilon_4 + \varepsilon_5.
\end{align*}
Then we have
\begin{align*}
	\Delta(\lie{p}_+ \cap (\lieC{h})^\perp, \lieC{t}) = \set{\varepsilon_1 + \varepsilon_2} \cup \set{\varepsilon_i + \varepsilon_j: 1 \leq i < 3 \leq j \leq 5}.
\end{align*}
Since $\set{\varepsilon_1 + \varepsilon_3, \varepsilon_2 + \varepsilon_4}$ is a set of strongly orthogonal roots in $\Delta(\lie{p}_+ \cap (\lieC{h})^\perp, \lieC{t})$, we have $r = 2 = \rank_{\RR}(\lie{g}^{\theta\sigma})$.
The assumptions of Proposition \ref{prop:NonMinimalFullRank} hold.

\subsubsection{\texorpdfstring{$\lie{g} = \lieE_{7(7)}$}{g=e7(7)}}

% split

We have $\lie{g}^{\theta\sigma} \simeq \lieSU(4, 4)$.
Then $r$ in Subsection \ref{subsect:AssociatedPair} satisfies
\begin{align*}
	r \leq \rank_{\RR}(\lie{g}^{\theta\sigma}) = 4 < 6 = \rank(\lie{g}) - 1.
\end{align*}
There are two embeddings of $\lieSU(1, k)$ into $\lieE_{7(7)}$.

\begin{align*}
	(\lie{g}, \lieNorm_{\lie{g}}(\lie{h})): \quad & \dynkin[scale=2, extended, labels={\alpha_1, \alpha_2, \alpha_5, \alpha_3, \alpha_4, , \beta_1, \beta_2}, affine mark=*]{E}{oooo*oo} \\
	(\lie{g}^{\theta\sigma}, \lieNorm_{\lie{g}}(\lie{h})^{\theta\sigma}): \quad & \dynkin[scale=2, labels={\alpha_5, \alpha_4, \alpha_3, \alpha_2 + \alpha_1, , \beta_2, \beta_1}]{A}{ooo*ooo}
\end{align*}

In the first embedding, we have $(\lie{g}^{\theta\sigma}, \lieNorm_{\lie{g}}(\lie{h})^{\theta\sigma}) \simeq (\lieSU(4, 4), \lie{s}(\lieU(4, 1) \oplus \lieU(3)))$, which is a symmetric pair.
The assumptions of Proposition \ref{prop:NonMinimalSymmetric} hold.


\begin{align*}
	(\lie{g}, \lieNorm_{\lie{g}}(\lie{h})): \quad & \dynkin[scale=2, extended, labels={\alpha_1, \alpha_2, , \alpha_3, \alpha_4, , \beta_1, \beta_2}, affine mark=*]{E}{oooo*oo} \\
	(\lie{g}^{\theta\sigma}, \lieNorm_{\lie{g}}(\lie{h})^{\theta\sigma}): \quad & \dynkin[scale=2, labels={, \alpha_4, \alpha_3, \alpha_2 + \alpha_1, , \beta_2, \beta_1}]{A}{ooo*ooo}
\end{align*}

In the second embedding, we have $(\lie{g}^{\theta\sigma}, \lieNorm_{\lie{g}}(\lie{h})^{\theta\sigma}) \simeq (\lieSU(4, 4), \lie{s}(\lieU(1, 0) \oplus \lieU(3, 1) \oplus \lieU(0, 3)))$.
Taking the standard realization of the root system of $\lieSL(8, \CC)$, we have
\begin{align*}
	\Delta(\lieC{g}^{\theta\sigma}, \lieC{t}) &= \set{\pm (\varepsilon_i - \varepsilon_j): 1\leq i < j \leq 8}, \\
	\Delta(\lie{p}_+ \cap (\lieC{h})^\perp, \lieC{t}) &= \set{\varepsilon_i - \varepsilon_j: 1\leq i \leq 4,  6\leq j \leq 8} \\
	&\cup \set{\varepsilon_1 - \varepsilon_5}.
\end{align*}
Then $\set{\varepsilon_1 - \varepsilon_5, \varepsilon_2 - \varepsilon_6, \varepsilon_3 - \varepsilon_7, \varepsilon_4 - \varepsilon_8}$ is a set of strongly orthogonal roots in $\Delta(\lie{p}_+ \cap (\lieC{h})^\perp, \lieC{t})$.
Hence we have $r = 4 = \rank_{\RR}(\lie{g}^{\theta\sigma})$.
The assumptions of Proposition \ref{prop:NonMinimalFullRank} hold.

\subsubsection{\texorpdfstring{$\lie{g} = \lieE_{7(-5)}$}{g=e7(-5)}}

% quaternionic

We have $\lie{g}^{\theta\sigma} \simeq \lieSO^*(12) \oplus \lieSL(2,\RR)$.
Then $r$ in Subsection \ref{subsect:AssociatedPair} satisfies
\begin{align*}
	r \leq \rank_{\RR}(\lie{g}^{\theta\sigma}) = 4 < 6 = \rank(\lie{g}) - 1.
\end{align*}
There are three embeddings of $\lieSU(1, k)$ into $\lieE_{7(-5)}$.

\begin{align*}
	(\lie{g}, \lieNorm_{\lie{g}}(\lie{h})): \quad & \dynkin[scale=2, extended, labels={\alpha_1, \alpha_2, \beta_1, , \beta_2, \beta_3, \beta_4, \beta_5}, affine mark=*]{E}{oo*oooo} \\
	(\lie{g}^{\theta\sigma}, \lieNorm_{\lie{g}}(\lie{h})^{\theta\sigma}): \quad & \dynkin[scale=2, labels={\beta_1, \beta_2, \beta_3, \beta_4, \beta_5, }, label directions={,,,right}]{D}{ooooo*}
	\dynkin[scale=2, labels={\alpha_2 + \alpha_1}]{A}{*}
\end{align*}
In the first embedding, we have $(\lie{g}^{\theta\sigma}, \lieNorm_{\lie{g}}(\lie{h})^{\theta\sigma}) \simeq (\lieSO^*(12) \oplus \lieSL(2,\RR), \lieU(6) \oplus\lieSL(2,\RR))$, which is a symmetric pair.
The assumptions of Proposition \ref{prop:NonMinimalSymmetric} hold.

\begin{align*}
	(\lie{g}, \lieNorm_{\lie{g}}(\lie{h})): \quad & \dynkin[scale=2, extended, labels={\alpha_1, \alpha_2, , \alpha_3, \alpha_4, \alpha_5, \alpha_6, }, affine mark=*]{E}{o*oooo*} \\
	(\lie{g}^{\theta\sigma}, \lieNorm_{\lie{g}}(\lie{h})^{\theta\sigma}): \quad & \dynkin[scale=2, labels={\alpha_6, \alpha_5, \alpha_4, \alpha_3, , \alpha_2 + \alpha_1}, label directions={,,,right}]{D}{ooooo*}
	\dynkin[scale=2, labels={}]{A}{*}
\end{align*}
In the second embedding, we have $(\lie{g}^{\theta\sigma}, \lieNorm_{\lie{g}}(\lie{h})^{\theta\sigma}) \simeq (\lieSO^*(12) \oplus \lieSL(2,\RR), \lieU(5,1) \oplus\lieSO(2,\RR))$, which is a symmetric pair.
The assumptions of Proposition \ref{prop:NonMinimalSymmetric} hold.


\begin{align*}
	(\lie{g}, \lieNorm_{\lie{g}}(\lie{h})): \quad & \dynkin[scale=2, extended, labels={\alpha_1, \alpha_2, , \alpha_3, \alpha_4, , \beta_1, \beta_2}, affine mark=*]{E}{o*oo*oo} \\
	(\lie{g}^{\theta\sigma}, \lieNorm_{\lie{g}}(\lie{h})^{\theta\sigma}): \quad & \dynkin[scale=2, labels={\beta_1, \beta_2, , \alpha_3, \alpha_4, \alpha_2 + \alpha_1}, label directions={,,,right}]{D}{ooooo*}
	\dynkin[scale=2, labels={}]{A}{*}
\end{align*}

In the third embedding, we have $(\lie{g}^{\theta\sigma}, \lieNorm_{\lie{g}}(\lie{h})^{\theta\sigma}) \simeq (\lieSO^*(12) \oplus \lieSL(2,\RR), \lieSO^*(6) \oplus \lieU(3, 0) \oplus\lieSO(2))$.
Taking the standard realizations of the root systems of $\lieSO(12,\CC)$ and $\lieSL(2,\CC)$, we have
\begin{align*}
	\Delta(\lieC{g}^{\theta\sigma}, \lieC{t}) &= \set{\pm \varepsilon_i \pm \varepsilon_j: 1\leq i < j \leq 6} \cup \set{\pm \varepsilon_7}, \\
	\Delta(\lie{p}_+ \cap (\lieC{h})^\perp, \lieC{t}) &= \set{\varepsilon_i + \varepsilon_j: 1\leq i < j \leq 3} \\
	&\cup \set{\varepsilon_i + \varepsilon_j: 1\leq i \leq 3, 4\leq j \leq 6} \cup \set{\varepsilon_7}.
\end{align*}
Then $\set{\varepsilon_1 + \varepsilon_4, \varepsilon_2 + \varepsilon_5, \varepsilon_3 + \varepsilon_6, \varepsilon_7}$ is a set of strongly orthogonal roots in $\Delta(\lie{p}_+ \cap (\lieC{h})^\perp, \lieC{t})$.
Hence we have $r = 4 = \rank_{\RR}(\lie{g}^{\theta\sigma})$.
The assumptions of Proposition \ref{prop:NonMinimalFullRank} hold.

\subsubsection{\texorpdfstring{$\lie{g} = \lieE_{7(-25)}$}{g=e7(-25)}}

% Hermitian

We have $\lie{g}^{\theta\sigma} \simeq \lieE_{6(-14)} \oplus \lieSO(2)$.
Then $r$ in Subsection \ref{subsect:AssociatedPair} satisfies
\begin{align*}
	r \leq \rank_{\RR}(\lie{g}^{\theta\sigma}) = 2 < 6 = \rank(\lie{g}) - 1.
\end{align*}
There is one embedding of $\lieSU(1, k)$ into $\lieE_{7(-25)}$.

\begin{align*}
	(\lie{g}, \lieNorm_{\lie{g}}(\lie{h})): \quad & \dynkin[scale=2, extended, labels={\alpha_1, \alpha_2, , \alpha_3, \alpha_4, \alpha_5, \alpha_6, }, affine mark=*]{E}{oooooo*} \\
	(\lie{g}^{\theta\sigma}, \lieNorm_{\lie{g}}(\lie{h})^{\theta\sigma}): \quad & \dynkin[scale=2, labels={\alpha_2+\alpha_1, , \alpha_3, \alpha_4, \alpha_5, \alpha_6}]{E}{*ooooo}
\end{align*}

In the embedding, we have $(\lie{g}^{\theta\sigma}, \lieNorm_{\lie{g}}(\lie{h})^{\theta\sigma}) \simeq (\lieE_{6(-14)} \oplus \lieSO(2), \lieSU(1, 5) \oplus \lieSO(2) \oplus \lieSO(2))$.
Then the subalgebra $\lieNorm_{\lie{g}}(\lie{h})^{\theta\sigma}$ is contained in a symmetric subalgebra $(\lie{g}^{\theta\sigma})^{\tau}$ isomorphic to $\lieSU(1,5) \oplus \lieSL(2,\RR)\oplus \lieSO(2)$, where $\tau$ is an involution commuting with $\theta$.
By \cite[Table I]{OsSe84}, we have $(\lie{g}^{\theta\sigma})^{\theta\tau} \simeq \lieSO^*(10) \oplus \lieSO(2) \oplus \lieSO(2)$ and
\begin{align*}
	\Delta(\lie{p}_+ \cap (\lieC{h})^\perp, \lieC{t}) \supset \Delta(\lie{p}_+ \cap (\lieC{g}^{\theta\sigma})^{\theta \tau}, \lieC{t}).
\end{align*}
By the construction of a maximal set of strongly orthogonal roots (\cite[Lemma 7.143]{Kn02}), we have $2 = \rank_{\RR}((\lie{g}^{\theta\sigma})^{\theta\tau}) \leq r$.
Hence we obtain $r = 2 = \rank_{\RR}(\lie{g}^{\theta\sigma})$.
The assumptions of Proposition \ref{prop:NonMinimalFullRank} hold.

\subsubsection{\texorpdfstring{$\lie{g} = \lieE_{8(8)}$}{g=e8(8)}}

% split

We have $\lie{g}^{\theta\sigma} \simeq \lieSO^*(16)$.
Then $r$ in Subsection \ref{subsect:AssociatedPair} satisfies
\begin{align*}
	r \leq \rank_{\RR}(\lie{g}^{\theta\sigma}) = 4 < 7 = \rank(\lie{g}) - 1.
\end{align*}
There are four embeddings of $\lieSU(1, k)$ into $\lieE_{8(8)}$.

\begin{align*}
	(\lie{g}, \lieNorm_{\lie{g}}(\lie{h})): \quad & \dynkin[scale=2, extended, labels={\alpha_1, \alpha_8, , \alpha_7, \alpha_6, \alpha_5, \alpha_4, \alpha_3, \alpha_2}, affine mark=*]{E}{o*oooooo} \\
	(\lie{g}^{\theta\sigma}, \lieNorm_{\lie{g}}(\lie{h})^{\theta\sigma}): \quad & \dynkin[scale=2, labels={\alpha_8, \alpha_7, \alpha_6, \alpha_5, \alpha_4, \alpha_3, , \alpha_2+\alpha_1}, label directions={,,,,,right}]{D}{ooooooo*}
\end{align*}

In the first embedding, we have $(\lie{g}^{\theta\sigma}, \lieNorm_{\lie{g}}(\lie{h})^{\theta\sigma}) \simeq (\lieSO^*(16), \lieU(1, 7))$, which is a symmetric pair.
The assumptions of Proposition \ref{prop:NonMinimalSymmetric} hold.

\begin{align*}
	(\lie{g}, \lieNorm_{\lie{g}}(\lie{h})): \quad & \dynkin[scale=2, extended, labels={\alpha_1, , , \alpha_7, \alpha_6, \alpha_5, \alpha_4, \alpha_3, \alpha_2}, affine mark=*]{E}{o*oooooo} \\
	(\lie{g}^{\theta\sigma}, \lieNorm_{\lie{g}}(\lie{h})^{\theta\sigma}): \quad & \dynkin[scale=2, labels={, \alpha_7, \alpha_6, \alpha_5, \alpha_4, \alpha_3, , \alpha_2+\alpha_1}, label directions={,,,,,right}]{D}{ooooooo*}
\end{align*}

In the second embedding, we have $(\lie{g}^{\theta\sigma}, \lieNorm_{\lie{g}}(\lie{h})^{\theta\sigma}) \simeq (\lieSO^*(16), \lieU(1, 6) \oplus \lieU(1))$.
Taking the standard realization of the root system of $\lieSO(16, \CC)$, we have
\begin{align*}
	\Delta(\lieC{g}^{\theta\sigma}, \lieC{t}) &= \set{\pm \varepsilon_i \pm \varepsilon_j: 1\leq i < j \leq 8},\\
	\Delta(\lie{p}_+ \cap (\lieC{h})^\perp, \lieC{t}) &= \set{\varepsilon_i + \varepsilon_j: 1\leq i < j \leq 7} \\
	&\cup \set{\varepsilon_1 + \varepsilon_8}.
\end{align*}
Then $\set{\varepsilon_1 + \varepsilon_8, \varepsilon_2 + \varepsilon_7, \varepsilon_3 + \varepsilon_6, \varepsilon_4 + \varepsilon_5}$ is a set of strongly orthogonal roots in $\Delta(\lie{p}_+ \cap (\lieC{h})^\perp, \lieC{t})$.
Hence we have $r = 4 = \rank_{\RR}(\lie{g}^{\theta\sigma})$.
The assumptions of Proposition \ref{prop:NonMinimalFullRank} hold.

\begin{align*}
	(\lie{g}, \lieNorm_{\lie{g}}(\lie{h})): \quad & \dynkin[scale=2, extended, labels={\alpha_1, \beta_1, , , \alpha_6, \alpha_5, \alpha_4, \alpha_3, \alpha_2}, affine mark=*]{E}{o*oooooo} \\
	(\lie{g}^{\theta\sigma}, \lieNorm_{\lie{g}}(\lie{h})^{\theta\sigma}): \quad & \dynkin[scale=2, labels={\beta_1, , \alpha_6, \alpha_5, \alpha_4, \alpha_3, , \alpha_2+\alpha_1}, label directions={,,,,,right}]{D}{ooooooo*}
\end{align*}

In the third embedding, we have $(\lie{g}^{\theta\sigma}, \lieNorm_{\lie{g}}(\lie{h})^{\theta\sigma}) \simeq (\lieSO^*(16), \lieU(1, 5) \oplus \lieU(2))$.
Taking the standard realization of the root system of $\lieSO(16, \CC)$, we have
\begin{align*}
	\Delta(\lieC{g}^{\theta\sigma}, \lieC{t}) &= \set{\pm \varepsilon_i \pm \varepsilon_j: 1\leq i < j \leq 8},\\
	\Delta(\lie{p}_+ \cap (\lieC{h})^\perp, \lieC{t}) &= \set{\varepsilon_i + \varepsilon_j: 1\leq i < j \leq 7} \\
	&\cup \set{\varepsilon_1 + \varepsilon_8, \varepsilon_2 + \varepsilon_8}.
\end{align*}
Then $\set{\varepsilon_1 + \varepsilon_8, \varepsilon_2 + \varepsilon_7, \varepsilon_3 + \varepsilon_6, \varepsilon_4 + \varepsilon_5}$ is a set of strongly orthogonal roots in $\Delta(\lie{p}_+ \cap (\lieC{h})^\perp, \lieC{t})$.
Hence we have $r = 4 = \rank_{\RR}(\lie{g}^{\theta\sigma})$.
The assumptions of Proposition \ref{prop:NonMinimalFullRank} hold.

\begin{align*}
	(\lie{g}, \lieNorm_{\lie{g}}(\lie{h})): \quad & \dynkin[scale=2, extended, labels={\alpha_1, \beta_1, \beta_4, \beta_2, \beta_3, , \alpha_4, \alpha_3, \alpha_2}, affine mark=*]{E}{oooo*ooo} \\
	(\lie{g}^{\theta\sigma}, \lieNorm_{\lie{g}}(\lie{h})^{\theta\sigma}): \quad & \dynkin[scale=2, labels={\beta_1, \beta_2, \beta_3, \beta_4, , \alpha_3, \alpha_4, \alpha_2+\alpha_1}, label directions={,,,,,right}]{D}{ooooooo*}
\end{align*}

In the fourth embedding, we have $(\lie{g}^{\theta\sigma}, \lieNorm_{\lie{g}}(\lie{h})^{\theta\sigma}) \simeq (\lieSO^*(16), \lieU(5) \oplus \lieSO^*(6))$.
Taking the standard realization of the root system of $\lieSO(16, \CC)$, we have
\begin{align*}
	\Delta(\lieC{g}^{\theta\sigma}, \lieC{t}) &= \set{\pm \varepsilon_i \pm \varepsilon_j: 1\leq i < j \leq 8},\\
	\Delta(\lie{p}_+ \cap (\lieC{h})^\perp, \lieC{t}) &= \set{\varepsilon_i + \varepsilon_j: 1\leq i < j \leq 5} \\
	&\cup \set{\varepsilon_i + \varepsilon_j: 1\leq i \leq 5, 6\leq j \leq 8}.
\end{align*}
Then $\set{\varepsilon_1 + \varepsilon_8, \varepsilon_2 + \varepsilon_7, \varepsilon_3 + \varepsilon_6, \varepsilon_4 + \varepsilon_5}$ is a set of strongly orthogonal roots in $\Delta(\lie{p}_+ \cap (\lieC{h})^\perp, \lieC{t})$.
Hence we have $r = 4 = \rank_{\RR}(\lie{g}^{\theta\sigma})$.
The assumptions of Proposition \ref{prop:NonMinimalFullRank} hold.

\subsubsection{\texorpdfstring{$\lie{g} = \lieE_{8(-24)}$}{g=e8(-24)}}

% quaternionic

We have $\lie{g}^{\theta\sigma} \simeq \lieE_{7(-25)} \oplus \lieSL(2,\RR)$.
Then $r$ in Subsection \ref{subsect:AssociatedPair} satisfies
\begin{align*}
	r \leq \rank_{\RR}(\lie{g}^{\theta\sigma}) = 4 < 7 = \rank(\lie{g}) - 1.
\end{align*}
The Vogan diagram of $\lie{g}^{\theta\sigma}$ is given by
\begin{align*}
	\dynkin[scale=2, labels={\delta, \gamma_1, \gamma_2, \gamma_3, \gamma_4, \gamma_5, \gamma_6}]{E}{oooooo*}
	\quad \dynkin[scale=2, labels={\varepsilon_0}]{A}{*}
\end{align*}
Here the labels are given by a realization of the root system $\Delta(\lieC{g}^{\theta\sigma}, \lieC{t})$ in the Euclidean space $\bigoplus_{i=0}^8 \RR \varepsilon_i$ such that
\begin{align*}
	\Delta^+(\lieC{g}^{\theta\sigma}, \lieC{t}) &= \set{\varepsilon_0} \cup \set{\varepsilon_j \pm \varepsilon_i: 1\leq i < j \leq 6} \cup \set{\varepsilon_8 - \varepsilon_7}\\
	&\cup \set{\frac{1}{2}\sum_{i=1}^8 a_i \varepsilon_i: a_i = \pm 1, a_7=-1, a_8=1, a_1a_2\cdots a_6 = -1}, \\
	\Delta^+(\lieC{g}^{\theta, \sigma}, \lieC{t}) &= \set{\varepsilon_j \pm \varepsilon_i: 1\leq i < j \leq 5} \\
	& \cup \set{\frac{1}{2}\sum_{i=1}^8 a_i \varepsilon_i: a_i = \pm 1, a_6=a_7=-1,a_8=1,a_1a_2\cdots a_5 = 1},\\
	\gamma_i &= \varepsilon_i - \varepsilon_{i - 1} \quad (2\leq i \leq 6), \quad \gamma_1 = \varepsilon_2 + \varepsilon_1, \\
	\delta &= \frac{1}{2}(\varepsilon_8 - \varepsilon_7 - \varepsilon_6 - \varepsilon_5 - \varepsilon_4 - \varepsilon_3 - \varepsilon_2 + \varepsilon_1)
\end{align*}
See \cite[Chapter VI, \S 4.11]{Bo02}.
Then we have
\begin{align*}
	\Delta(\lie{p}_+, \lieC{t}) &= \set{\varepsilon_0} \cup \set{\varepsilon_6 \pm \varepsilon_i: 1\leq i \leq 5} \cup \set{\varepsilon_8 - \varepsilon_7} \\
	&\cup \set{\frac{1}{2}\sum_{i=1}^8 a_i \varepsilon_i: a_i = \pm 1, a_7=-1,a_6=a_8=1,a_1a_2\cdots a_5 = -1}.
\end{align*}

There are three embeddings of $\lieSU(1, k)$ into $\lieE_{8(-24)}$.

\begin{align*}
	(\lie{g}, \lieNorm_{\lie{g}}(\lie{h})): \quad & \dynkin[scale=2, extended, labels={\alpha_1, \beta_1, \beta_2, \beta_3, \beta_4, \beta_5, \beta_6, , \alpha_2}, affine mark=*]{E}{oooooo*o} \\
	(\lie{g}^{\theta\sigma}, \lieNorm_{\lie{g}}(\lie{h})^{\theta\sigma}): \quad & \dynkin[scale=2, labels={\beta_1, \beta_2, \beta_3, \beta_4, \beta_5, \beta_6, }]{E}{oooooo*}
	\dynkin[scale=2, labels={\alpha_2 + \alpha_1}]{A}{*}
\end{align*}

In the first embedding, we have $(\lie{g}^{\theta\sigma}, \lieNorm_{\lie{g}}(\lie{h})^{\theta\sigma}) \simeq (\lieE_{7(-25)} \oplus \lieSL(2,\RR), \lieE_{6(-78)} \oplus \lieSO(2) \oplus \lieSL(2,\RR))$, which is a symmetric pair.
The assumptions of Proposition \ref{prop:NonMinimalSymmetric} hold.

\begin{align*}
	(\lie{g}, \lieNorm_{\lie{g}}(\lie{h})): \quad & \dynkin[scale=2, extended, labels={\alpha_1, \beta_1, , , \alpha_6, \alpha_5, \alpha_4, \alpha_3, \alpha_2}, affine mark=*]{E}{o**ooooo} \\
	(\lie{g}^{\theta\sigma}, \lieNorm_{\lie{g}}(\lie{h})^{\theta\sigma}): \quad & \dynkin[scale=2, labels={\beta_1, \alpha_6, , \alpha_5, \alpha_4, \alpha_3, \alpha_2 + \alpha_1}]{E}{oooooo*}
	\quad \dynkin[scale=2, labels={}]{A}{*}
\end{align*}

In the second embedding, we have $(\lie{g}^{\theta\sigma}, \lieNorm_{\lie{g}}(\lie{h})^{\theta\sigma}) \simeq (\lieE_{7(-25)} \oplus \lieSL(2,\RR), \lieSU(1, 5) \oplus \lieU(1) \oplus \lieSU(2) \oplus \lieSO(2))$.
Then we have
\begin{align*}
	&\Delta(\lie{p}_+ \cap (\lieC{h})^\perp, \lieC{t}) \\
	&= \set{\varepsilon_0} \cup \set{\varepsilon_6 + \varepsilon_i: 2\leq i \leq 5} \cup \set{\varepsilon_6 - \varepsilon_1} \cup \set{\varepsilon_8 - \varepsilon_7} \\
	&\cup \set{\frac{1}{2}\sum_{i=1}^8 a_i \varepsilon_i: a_i = \pm 1, a_7=-1,a_6=a_8=1,a_1a_2\cdots a_5 = -1}.
\end{align*}
Hence
\begin{align*}
	\set{\begin{aligned}
		&\varepsilon_0, \varepsilon_6 + \varepsilon_5, \\
		&\frac{1}{2}(\varepsilon_1 + \varepsilon_2 - \varepsilon_3 - \varepsilon_4 - \varepsilon_5 + \varepsilon_6 - \varepsilon_7 + \varepsilon_8), \\
		&\frac{1}{2}(-\varepsilon_1 - \varepsilon_2 + \varepsilon_3 + \varepsilon_4 - \varepsilon_5 + \varepsilon_6 - \varepsilon_7 + \varepsilon_8)
	\end{aligned}}
\end{align*}
is a set of strongly orthogonal roots in $\Delta(\lie{p}_+ \cap (\lieC{h})^\perp, \lieC{t})$.
Hence we have $r = 4 = \rank_{\RR}(\lie{g}^{\theta\sigma})$.
The assumptions of Proposition \ref{prop:NonMinimalFullRank} hold.

\begin{align*}
	(\lie{g}, \lieNorm_{\lie{g}}(\lie{h})): \quad & \dynkin[scale=2, extended, labels={\alpha_1, , , \alpha_7, \alpha_6, \alpha_5, \alpha_4, \alpha_3, \alpha_2}, affine mark=*]{E}{**oooooo} \\
	(\lie{g}^{\theta\sigma}, \lieNorm_{\lie{g}}(\lie{h})^{\theta\sigma}): \quad & \dynkin[scale=2, labels={\alpha_7, , \alpha_6, \alpha_5, \alpha_4, \alpha_3, \alpha_2 + \alpha_1}]{E}{oooooo*}
	\quad \dynkin[scale=2, labels={}]{A}{*}
\end{align*}
In the third embedding, we have $(\lie{g}^{\theta\sigma}, \lieNorm_{\lie{g}}(\lie{h})^{\theta\sigma}) \simeq (\lieE_{7(-25)} \oplus \lieSL(2,\RR), \lieU(1, 6) \oplus \lieSO(2))$.
Then we have
\begin{align*}
	&\Delta(\lie{p}_+ \cap (\lieC{h})^\perp, \lieC{t}) \\
	&= \set{\varepsilon_0} \cup \set{\varepsilon_6 + \varepsilon_i: 1\leq i \leq 5} \cup \set{\varepsilon_8 - \varepsilon_7} \\
	&\cup \set{\frac{1}{2}\sum_{i=1}^8 a_i \varepsilon_i: a_i = \pm 1, a_7=-1,a_6=a_8=1,a_1a_2\cdots a_5 = -1}\\
	&\backslash \set{\frac{1}{2}\left(-\sum_{i=1}^5 \varepsilon_i + \varepsilon_6 - \varepsilon_7 + \varepsilon_8\right)}.
\end{align*}
Hence
\begin{align*}
	\set{\begin{aligned}
		&\varepsilon_0, \varepsilon_6 + \varepsilon_5,\\
		&\frac{1}{2}(\varepsilon_1 + \varepsilon_2 - \varepsilon_3 - \varepsilon_4 - \varepsilon_5 + \varepsilon_6 - \varepsilon_7 + \varepsilon_8), \\
		&\frac{1}{2}(-\varepsilon_1 - \varepsilon_2 + \varepsilon_3 + \varepsilon_4 - \varepsilon_5 + \varepsilon_6 - \varepsilon_7 + \varepsilon_8)
	\end{aligned}}
\end{align*}
is a set of strongly orthogonal roots in $\Delta(\lie{p}_+ \cap (\lieC{h})^\perp, \lieC{t})$.
Hence we have $r = 4 = \rank_{\RR}(\lie{g}^{\theta\sigma})$.
The assumptions of Proposition \ref{prop:NonMinimalFullRank} hold.

\subsubsection{\texorpdfstring{$\lie{g} = \lieF_{4(4)}$}{g=F4(4)}}

% split
For $\lie{g} = \lieF_{4(4)}$, we check the assumptions of Proposition \ref{prop:ReductionRealStronglyOrthogonal}.

\begin{align*}
	(\lie{g}, \lieNorm_{\lie{g}}(\lie{h})): \quad & \dynkin[scale=2, extended, labels={\alpha_1, \alpha_2, , \beta_1, \beta_2}, affine mark=*]{F}{o*oo} \\
	(\lie{g}^{\theta\sigma}, \lieNorm_{\lie{g}}(\lie{h})^{\theta\sigma}): \quad & \dynkin[scale=2, labels={\beta_2, \beta_1}]{C}{oo*}
	\dynkin[scale=2, labels={\alpha_2 + \alpha_1}]{A}{*}
\end{align*}

In the embedding, we have $(\lie{g}^{\theta\sigma}, \lieNorm_{\lie{g}}(\lie{h})^{\theta\sigma}) \simeq (\lieSp(3, \RR) \oplus \lieSL(2,\RR), \lieU(3) \oplus \lieSL(2, \RR))$, which is a symmetric pair.
Let $\lie{g}^{\theta\sigma} = \lie{i} \oplus \lie{j}$ be the decomposition into simple ideals such that $\lie{j} \simeq \lieSp(3, \RR)$.

\begin{lemma} \label{lem:ProjectionToGPrime}
	The projection of $\lie{g}_0\cap \lie{k}$ to $\lie{j}$ coincides with $\lie{j}\cap \lie{k}$.
\end{lemma}

\begin{proof}
	Let $Z\in \lie{g}_2$ be the element such that $\ad(Z+\theta(Z))$ defines a complex structure on $\lie{g}_1 \oplus \lie{g}_{-1}$ as in Lemma \ref{lem:ComplexStructure}.
	Then $\ad(Z+\theta(Z))$ gives a complex structure on $\lie{g}^{-\sigma, -\theta}$.
	Hence $Z+\theta(Z)$ has non-zero projections to both factors $\lie{i}$ and $\lie{j}$.
	Since $(\lie{g}_0\cap \lie{k}) \oplus \RR (Z+\theta(Z)) = \lie{k}^{\sigma}$ is an orthogonal direct sum, the projection of $\lie{g}_0\cap \lie{k}$ to $\lie{j}$ coincides with $\lie{j}\cap \lie{k}$.
\end{proof}

By the isomorphism $(\lie{g}^{\theta\sigma}, \lieNorm_{\lie{g}}(\lie{h})^{\theta\sigma}) \simeq (\lieSp(3, \RR) \oplus \lieSL(2,\RR), \lieU(3) \oplus \lieSL(2, \RR))$ and Lemma \ref{lem:ProjectionToGPrime}, the assumptions of Proposition \ref{prop:ReductionRealStronglyOrthogonal} hold.


\subsubsection{\texorpdfstring{$\lie{g} = \lieG_{2(2)}$}{g=G2(2)}}

% split

For $\lie{g} = \lieG_{2(2)}$, we check the assumptions of Proposition \ref{prop:ReductionRealStronglyOrthogonal}.

\begin{align*}
	(\lie{g}, \lieNorm_{\lie{g}}(\lie{h})): \quad & \dynkin[scale=2, extended, labels={\alpha_1, \alpha_2, }, affine mark=*]{G}{o*} \\
	(\lie{g}^{\theta\sigma}, \lieNorm_{\lie{g}}(\lie{h})^{\theta\sigma}): \quad & \dynkin[scale=2, labels={\alpha_2+\alpha_1}]{A}{*}
	\dynkin[scale=2, labels={}]{A}{*}
\end{align*}

In the embedding, we have $(\lie{g}^{\theta\sigma}, \lieNorm_{\lie{g}}(\lie{h})^{\theta\sigma}) \simeq (\lieSL(2, \RR) \oplus \lieSL(2,\RR), \lieSL(2, \RR) \oplus \lieSO(2))$, which is a symmetric pair.
Let $\lie{g}^{\theta\sigma} = \lie{i} \oplus \lie{j}$ be the decomposition into simple ideals such that $\lie{i}\subset \lie{h}$.
By the same argument as in the proof of Lemma \ref{lem:ProjectionToGPrime}, the projection of $\lie{g}_0\cap \lie{k}$ to $\lie{j}$ coincides with $\lie{j}\cap \lie{k}$.
Therefore, by the isomorphism $(\lie{g}^{\theta\sigma}, \lieNorm_{\lie{g}}(\lie{h})^{\theta\sigma}) \simeq (\lieSL(2, \RR) \oplus \lieSL(2,\RR), \lieSL(2, \RR) \oplus \lieSO(2))$, the assumptions of Proposition \ref{prop:ReductionRealStronglyOrthogonal} hold.

We have completed the proof of Theorem \ref{thm:NonMinimalExceptional}, and therefore the proof of Theorem \ref{thm:MainTargetMinimal} in all cases.
